% ============================================================
%  H. Høffding, "Kontinuiteten i Kants filosofiske Udviklingsgang"
%  D. Kgl. Danske Vidensk. Selsk. Skr., 6. Række,
%  historisk og filosofisk Afd. IV. 1. Kjøbenhavn 1893.
%  TRANSCRIPTION (Danish)
%  Source: kants-udviklingsgang.pdf (Bianco Lunos, 1893).
%  Footnotes numbered continuously; in the original they
%  restart on each page.
% ============================================================
\documentclass[12pt,a4paper]{article}
\usepackage[utf8]{inputenc}
\usepackage[T1]{fontenc}
\usepackage{libertinus}
\usepackage{libertinust1math}
\usepackage[danish]{babel}
\usepackage[protrusion=true,expansion=true]{microtype}
\usepackage{setspace}
\usepackage[margin=1.2in]{geometry}
\usepackage{fancyhdr}
\usepackage{hyperref}

\hypersetup{
  pdftitle={Kontinuiteten i Kants filosofiske Udviklingsgang},
  pdfauthor={Harald Høffding},
  hidelinks
}

\pagestyle{fancy}
\fancyhf{}
\rhead{Høffding, \emph{Kontinuiteten i Kants filosofiske Udviklingsgang} (1893)}
\cfoot{\thepage}

\setstretch{1.3}

\title{\textbf{Kontinuiteten i Kants filosofiske Udviklingsgang}}
\author{Harald Høffding}
\date{\emph{D.\ Kgl.\ Danske Vidensk.\ Selsk.\ Skr.}, 6.\ Række,\\
  historisk og filosofisk Afd.\ IV.\ 1.\ Kjøbenhavn 1893}

\begin{document}
\maketitle
\thispagestyle{fancy}

\section*{Indledning.}

\noindent 1.\ Forud for Kants Hovedværk «Kritik der reinen Vernunft»
(1781) gaar der som bekendt en lang Aarrække, i hvilken de store Ideer,
der gennem dette Værk skulde faa saa stor Indflydelse paa det følgende
Aarhundredes Aandsliv, langsomt modnedes. Kant er derfor ofte bleven
anført som Exempel paa sen Modenhed. Han var 57 Aar, da «Kritik der
reinen Vernunft» udkom. Men skønt hans Betydning og Berømthed utvivlsomt
især skyldes dette Værk, vilde det dog være urigtigt at antage, at hans
foregaaende Arbejder blot skulde have Værdi som forberedende Led i hans
Udviklingshistorie og miste deres Betydning, efter at Hovedværket var
kommet for Lyset. Som de allerede i sin Tid skaffede Kant et højt anset
Navn, saa at han endog formedelst dem blev stillet sammen med Lessing,
førend han stod som Forfatteren til «Kritik der reinen Vernunft»,
saaledes frembyde de endnu den Dag i Dag mere end rent historisk
Interesse. De indeholde Tanker, hvis Tid ikke er forbi. Mange af disse
Tanker har Kant --- til Dels med Metamorfoser, der betingedes ved hans
senere Standpunkt --- stedse fastholdt og paa ny udviklet i sine senere
Værker; det er især om dem, at denne Afhandling skal dreje sig. Men der
findes i de tidligere Arbejder ogsaa Tanker, som i de senere Værker med
Urette stilledes i Skygge eller helt opgaves. Det var ikke i alle
Henseender et Fremskridt, Kant gjorde, da han naaede sit definitive
Standpunkt. Særligt hvad Formen angaar, have de tidligere Arbejder et
afgjort Fortrin. Kant fremtræder i dem som en klar og elegant Forfatter,
hvad han ingenlunde kan siges at være i alle Partier af sine senere
Værker.

Det var naturligt, at det Nye, Storartede og Betydningsfulde i
Hovedværket baade for Kant selv, hans Disciple og hans Fortolkere maatte
faa de tidligere Arbejder til at træde i Skygge. Først i de senere Aar
har man af rent historisk Interesse draget dem frem. Men det har da
netop overvejende været for at paavise Overgangsleddene til det
definitive Standpunkt. Dette Studium har --- i Forbindelse med de
Kantiske Manuskripter, som i de senere Aar for første Gang ere bragte
for Lyset --- bragt Klarhed over Kants Udviklingsgang i flere
Henseender, samtidigt med, at det paa nogle Punkter har vist sig umuligt
at komme til noget bestemt Resultat. Jeg vil i denne Afhandling søge at
belyse nogle Hovedpunkter i Kants filosofiske Udvikling, saaledes at jeg
særligt vil godtgøre Kontinuiteten i denne og den blivende Betydning af
Adskilligt i de tidligere Arbejder.

Kant selv har ved forskellige Lejligheder udtalt, at hans Udvikling har
havt bestemte Vendepunkter. I et Brev til Mendelssohn af 18.\ August
1783 skriver han om den Maade, paa hvilken «Kritik der reinen Vernunft»
var bleven udarbejdet: «Hvad der var Produktet af i det Mindste tolv
Aars Eftertanke, havde jeg affattet i Løbet af fire til fem
Maaneder»\footnote{En aldeles lignende Udtalelse findes allerede i Kants
Brev til Garve af 7.\ August 1783: «Den Vortrag der Materien, die ich
mehr als 12 Jahre hinter einander sorgfältig durchgedacht hatte, \ldots\
brachte ich in etwa 4 bis 5 Monathen zu Stande.» (Dette Brev er trykt
hos Alb.\ Stern: \textit{Ueber die Beziehungen Chr.\ Garve's zu Kant},
p.\ 33 ff.)}. Denne Udtalelse viser os hen til Aaret 1769 som Tidspunktet
for Indledelsen af den Tankerække, der direkte førte til «Kritik der
reinen Vernunft». Og hermed stemmer hvad Kant langt tidligere havde
skrevet til Lambert (2.\ September 1770): «For et Aar siden er jeg, som
jeg smigrer mig med, naat til det Begreb, som jeg ikke frygter for
nogensinde at behøve at ændre, om jeg end maa udvide det, og ved hvilket
alle Slags metafysiske Spørgsmaal kunne prøves efter ganske sikre og
lette Kriterier, saa at det med Sikkerhed kan afgøres, hvor vidt de ere
opløselige eller ikke.» I de Optegnelser af Kant, som Benno Erdmann har
udgivet under Titelen «Kants Reflexioner», og som indeholde adskillige
interessante Bidrag til hans Udviklingshistorie, hedder det, i et
Fragment, der synes at høre til et Udkast til en paatænkt Fortale til
«Kritik der reinen Vernunft»: «Aaret 1769 gav mig meget
Lys»\footnote{\textit{Reflexionen Kants zur kritischen Philosophie. Aus
Kants handschriftlichen Aufzeichnungen herausgegeben von Benno Erdmann.}
II.\ Leipzig 1884.\ p.\ 4.\ (Nr.\ 4.)}. Og da en af Kants Disciple vilde
foranstalte en Samling af hans mindre Skrifter, ønskede Kant ikke, at
der blev optaget Arbejder deri, som vare ældre end 1770 (Brev til
Tieftrunk af 6.\ Februar 1798).

Her haves altsaa et Datum for Begyndelsen af den sidste Periode i Kants
Forsken, og et Bevis for, at han selv opfattede den som en sammenhængende
Helhed. Et Vendepunkt i hans Udviklingsgang er fastslaaet, og dermed
stilles den Opgave --- særligt naar man vil vise Udviklingsgangens
Kontinuitet --- at forklare, hvorledes dette Vendepunkt forberedes ved
og hænger sammen med den forudgaaende Periode.

Endnu et Problem har Kant selv stillet dem, der beskæftige sig med hans
Udviklingsgang, ved sin Henvisning til den Gæld, han stod i til David
Hume. «Jeg tilstaar frit», hedder det i Fortalen til «Prolegomena», «det
var Erindringen om David Hume, der for mange Aar siden afbrød min
dogmatiske Slummer og gav mine Undersøgelser paa den spekulative
Filosofis Omraade en helt anden Retning»\footnote{\textit{Prolegomena zu
jeder künftigen Metaphysik.} Riga 1783.\ p.\ 13.}. Her henvises der
altsaa ogsaa til en Afbrydelse, idet en Impuls ude fra griber ind og
medfører en helt ny Retning i Tankearbejdet. Spørgsmaalet bliver da
først og fremmest, paa hvilket Punkt i Kants Udvikling denne Vækkelse
skal sættes. Dette Spørgsmaal har, trods al den Lærdom og Skarpsindighed,
der er anvendt derpaa, ikke naat nogen enstemmig Besvarelse af
Kantforskerne. Der er intet Punkt, hvor der er et saadant Spring i Kants
Tankeudvikling, at en bestemt Indgriben af en fremmed Tankegang skulde
kunne siges at være en uundgaaelig Antagelse. Overalt kan Fremgangen
meget vel forklares som Fortsættelse af den forudgaaende Tids Arbejde. En
af de mest fremragende Kantforskere har antaget, at Aaret 1769, i hvilket
efter Kants udtrykkelige Erklæring hans definitive Filosofis Grundtanke
opstod, ogsaa kunde være det Tidspunkt, i hvilket Vækkelsen af den
dogmatiske Slummer blev til\footnote{Friederich Paulsen: \textit{Versuch
einer Entwickelungsgeschichte der Kantischen Erkenntnisstheorie.}
Leipzig 1875.\ p.\ 126 ff.}. Det vilde jo unægteligt være det Simpleste,
om de to af Kant selv omtalte Vendepunkter faldt sammen. Om det forholder
sig saaledes, kan først afgøres ved en nærmere Undersøgelse af den
Ændring, der skete i 1769. I og for sig er der ingen Rimelighed for, at
de to Tidspunkter falde sammen. Et er jo en Vækkelse, et Andet en
Opdagelse; Et er en ny Retning, et Andet et nyt Resultat. Hvad der skete
i Aaret 1769, kan have været et Fremskridt, der kun blev muligt ved, at
Vækkelsen var gaaet forud, maaske lang Tid forud. Det Lys, der i det Aar
opgik for Kant, kan han have maattet søge længe, før han fandt det. Og
dette bekræftes dels ved, at han i Brevene til Garve og Mendelssohn taler
om «mere end tolv Aars Arbejde», «i det Mindste tolv Aars Arbejde», dels
ved, at han i «Prolegomena» udtaler, at den Forskel, han [i
«Dissertationen» (1770)] gjorde mellem Sansningens og Forstandens
Elementarbegreber, først lykkedes «efter lang
Eftertanke»\footnote{\textit{Prolegomena zu jeder künftigen Metaphysik.}
Riga 1783.\ p.\ 119.}.

Det vilde ganske vist være af ikke ringe Interesse, om Vækkelsens
Tidspunkt bestemt kunde angives. Forholdet mellem Kant og Hume er et
Hovedpunkt i hele den nyere Filosofis Historie. Man vil kunne fremstille
denne Historie som en stor Dialog, i hvilken Humes Replik og Kants Svar
paa den staa som de mest betydningsfulde Indlæg. Baade ved Modsætningen
mellem dem og ved Slægtskabet mellem dem er Studiet af deres Værker en
Kilde til stedse ny Orientering med Hensyn til filosofiske
Principspørgsmaal. Det vilde da være af stor Interesse at faa at vide,
paa hvilket Punkt i sin Udvikling Kant egentligt ret hørte Humes Replik.
At det ikke var det Tidspunkt, da han første Gang læste Hume, ligger i
hans Udtryk: «Erindringen om David Hume». --- Det vilde tillige være af
almindelig psykologisk Interesse, om dette Spørgsmaal kunde belyses
nærmere; vi vilde derved faa et klassisk Exempel paa, hvorledes en
afgørende Paavirkning kan modtages paa eminent selvstændig Maade.

Men netop i de Vanskeligheder, der ere forbundne med at paavise
Vækkelsens Tidspunkt ser jeg et Bevis paa Kontinuiteten i Kants
Udvikling. Netop fordi Kants Tænkning var i stadig Udfoldelse, --- netop
paa Grund af hans stadige Stræben og utrættelige Selvkritik (i hvert
Tilfælde indtil det for ham definitive Standpunkt var naat), er det
vanskeligt at afgøre, naar og hvorledes de ydre Impulser, som paa enkelte
Punkter vides at have været bestemmende, have grebet ind. --- Disse
Bemærkninger gælde iøvrigt ikke blot Indflydelsen fra David Hume, men
ogsaa den Paavirkning, Kant har modtaget fra en anden stor Samtidig, og
som paa det moralfilosofiske Omraade greb ikke mindre dybt ind i Kants
Udvikling end Humes gjorde paa det erkendelsesteoretiske, --- nemlig den,
som kom fra Jean Jacques Rousseau. Heller ikke denne vilde man være
bleven opmærksom paa, hvis der ikke tilfældigvis var blevet opbevaret
Ytringer af Kant i den Henseende.

Den Vanskelighed, som her har vist sig, kommer ofte igen ved
Undersøgelsen af en Udviklingsproces, især paa det aandelige Omraade. At
bestemme det nøjagtige Forhold mellem indre Udfoldelse og ydre
Paavirkning, hvor der fra først af er store Kræfter i Bevægelse, kan let
blive en uløselig Opgave. Det kan her ofte synes, som om den indre
Udfoldelse kunde gøre Alt, naar den blot fik gunstige Betingelser, uden
at nogen særegen ydre Impuls behøvedes. Saaledes gaar det os netop ved
Humes og Rousseaus Indflydelse paa Kant; det er kun de ydre Vidnesbyrd,
der føre til at undersøge den. Hvor Kantforskerne ere uenige med
hverandre, anvende de da ogsaa i Debatten den Metode at vise, hvorledes
Fremskridtet paa det bestemte Tidspunkt, der af Modparten antoges for
Vækkelsens Tidspunkt, kan forstaas uden ydre Paavirkning. Denne Metode
kan saa meget des lettere gennemføres, som en Paavirkning ved en
betydelig Aands Udvikling meget ofte vil øve sin Indflydelse paa
indirekte Vis, idet den mere vil bestaa i Udløsning af den egne, hidtil
bundne Kraft end i en Meddelelse af Stof. Hvad der bliver Vækkelsens
Resultat, behøver da ikke at ligne den vækkende Paavirkning. Det Ord,
hvorved man vækkes, har jo ofte slet ikke Noget at gøre med den Gerning,
man vækkes til.

Skønt jeg i det Følgende lejlighedsvis tillader mig en Formodning om
Tidspunktet for Vækkelsen, tillægger jeg dog Striden derom størst
Interesse derved, at den har styrket Indtrykket af Kontinuiteten i Kants
Udvikling. Kant selv var iøvrigt heller ikke af den Mening, at hans
førkritiske Skrifter ganske skulde have mistet deres Betydning. I et
Udkast til Fortale til «Kritik der reinen Vernunft» siger han: «Ved denne
min Afhandling er mine tidligere metafysiske Skrifters Værdi aldeles
tilintetgjort. Kun Ideens Rigtighed vil jeg endnu forsøge at
redde»\footnote{\textit{Reflexionen.} II.\ p.\ 5 (Nr.\ 7).}. Der
anerkendes dog altsaa et Tankeslægtskab trods Vækkelsen og Vendepunktet.
Det er det, der i nærværende Afhandling særligt skal behandles.

For Oversigtens Skyld deler jeg Fremstillingen i fire Afsnit, saaledes at
vi inden for hvert af disse undersøge Forholdet mellem det tidligere og
det senere Standpunkt for et enkelt bestemt Punkts Vedkommende.

\section*{I.\\ Aarsagsbegrebet.}

\noindent 2.\ I Kants Studieaar optoges hans Interesse af Wolfs Filosofi
og Newtons Fysik. Til begge indtog han en selvstændig Stilling, hvilket
allerede af den Grund var en Nødvendighed, at der var en Modstrid mellem
dem, da Wolf endnu fulgte den kartesiansk-leibnizianske Fysik og bestred
Newtons Attraktionslære. Allerede denne Modstrid maatte vække
Eftertanken, og vi se da ogsaa allerede i Kants første selvstændige
Periode (1755--62)\footnote{Forud for denne Periode ligger kun en
Afhandling om Vurderingen af levende Kraft (1747) og to
naturvidenskabelige Smaaafhandlinger (1754).} Forsøg af ham paa at ændre
sit filosofiske Grundlag saaledes, at Newtons Fysik kunde komme til sin
Ret. Kant stod med sin Anerkendelse af Newton ikke ene i den Wolfiske
Skole\footnote{Kants Lærer, Martin Knutzen, havde allerede søgt at
forbinde Newtons Fysik med Wolfs Filosofi, og han indførte sin Discipel i
Studiet af begge. Smign.\ Benno Erdmann: \textit{Martin Knutzen und seine
Zeit. Ein Beitrag zur Geschichte der Wolfischen Schule und insbesondere
zur Entwickelungsgeschichte Kants.} Leipzig 1876. --- I Danmark erklærede
Wolfianeren Jens Kraft, Professor i Sorø, sig i sin «Cosmologie»
(København 1752) mod den kartesianske Fysik og for Newtons.}. Men han
optraadte paa en selvstændigere Maade end andre Wolfianere, og det ikke
blot overfor Wolf, men ogsaa overfor Newton. I sit berømte
Ungdomsskrift «Allgemeine Naturgeschichte und Theorie des Himmels»
(1755) tager han Problemet om den naturlige Sammenhæng i Universet op paa
det Punkt, hvor Newton ikke blot havde forladt det, men endog erklæret
det for uløseligt. Trods sin store Ærbødighed for Newton kunde Kant ikke
slaa sig til Ro ved hans Udtalelse om, at der ikke kunde gives nogen
naturlig Aarsag til, at Planeterne bevæge sig omkring Solen i koncentriske
Baner, i samme Retning og omtrent i samme Plan, og at et tilsvarende
Forhold finder Sted hvad Biplaneterne angaar. Kant ansaa Appel til en
overnaturlig Aarsag for uvidenskabelig og forsøgte at vise, at den samme
lovmæssige Orden, der nu forener alle Legemer og Elementer til et stort
Hele, ogsaa havde lagt sig for Dagen ved Verdenssystemernes Oprindelse.
Derved førtes han til sin kosmogoniske Hypotese, i hvilken han søgte at
vise, at Solsystemets (og i Analogi dermed andre Verdenssystemers)
Oprindelse kan forstaas ved Virken af Kræfter, som den Dag i Dag ere i
Virksomhed. Det var Aktualitetsprincipet, det af Newton opstillede
Forlangende om \textit{vera causa}, som han anvendte paa det i Tiden
Fjerne, medens Newton kun havde anvendt det paa det i Rummet Fjerne.

Kants Iver for at gennemføre en mekanisk Naturopfattelse hænger nøje
sammen med en filosofisk Tankegang, der allerede paa den Tid udviklede
sig hos ham. Det var, mente han, en falsk Forudsætning, man gik ud fra,
naar man mente, at Naturen overladt til sig selv kun vilde frembringe
Uorden og Kaos: ikke ved et Tilfælde, men i Følge sine egne Love
frembringer Naturen det Hensigtsmæssige. Og netop det Faktum, at de
forskellige Elementer saaledes fra først til sidst virke sammen og
omspændes af en og samme Naturorden, vidner om, at de i deres Væsen og
Udspring ikke ere uafhængige af hverandre. Selve den mekaniske
Sammenhæng mellem Alt i Verden fører altsaa til Antagelsen af en
Enhedsgrund for Alt. To Opfattelser stiller Kant sig i Modsætning til,
--- paa den ene Side den absolute Atomistik, der antager indbyrdes
uafhængige Atomer, som kun ved et Tilfælde komme til at virke sammen, ---
paa den anden Side Fysikoteologien, der antager en ydre Indgriben af
Guddommen i Verdensmaskineriet. De tre Begreber: Element, Lov og Kraft
vilde han ikke paa udvortes Maade have skilt ad\footnote{Om den inderlige
Sammenhøren af Begreberne Element og Kraft handler særligt en Afhandling
af Kant fra det følgende Aar: \textit{Monadologia physica} (1756), i
hvilken Atomerne (o:\ de absolute Atomer) opfattes som Kraftpunkter, paa
lignende Maade som i vore Dage Fechner har opfattet dem i sin
«Atomenlehre».}. Denne Betragtning blev især af Betydning ved Problemet
om Oprindelsen til den Harmoni og Hensigtsmæssighed, der findes i Verden.
Kant klager over, at en næsten almindelig Fordom har indtaget de fleste
Forskere mod Naturens Evne til at frembringe noget Ordnet ved sine
almindelige Love. I hans Øjne er netop den Orden og Harmoni, der viser
sig i, hvad der frembringes efter Naturens Love, et Vidnesbyrd om, at
«alle Tings Væsen have et fælles Udspring i et Grundvæsen». Den fælles
Forbindelse tyder paa en fælles Afhængighed. Jo mere man lærer Naturen at
kende, des mere overbevises man om, at Tingene i Verden i deres inderste
Væsen ikke ere sondrede fra hverandre og fremmede for hverandre.

I dette sit Ungdomsskrift finder Kant altsaa Forbindelsen mellem
videnskabelig Erkendelse og religiøs Tro ikke ved at bryde gennem den
Sammenhæng, hin finder i Naturen, men netop ved at gaa ud fra denne
Sammenhæng. Og han opfatter ikke Teleologi og Mekanisme som Modsætninger,
men fordrer det Hensigtsmæssige forklaret som Frugten af en Udvikling
efter almindelige Naturlove.

3.\ Den samme Tanke, som Kant praktisk havde anvendt i sin astronomiske
Hypotese, udtalte han samtidigt i abstrakt Form i sit Habilitationsskrift
om Erkendelsens Principer\footnote{\textit{Principiorum primorum
cognitionis metaphysicæ nova dilucidatio.} (Sämtl.\ Werke, herausg.\ von
Rosenkranz und Schubert.\ I, p.\ 1--44.)}. Han kalder den her
Koexistentsens Princip. Hvis man --- saaledes er Kants Tankegang --- vilde
antage, at de enkelte Ting (Substantser) i Verden existerede hver for
sig, uafhængige af hverandre, vilde man ikke fra nogen enkelt af dem
kunne slutte sig til de andres Existents. Men saaledes er det ikke.
Formedelst den gensidige Sammenhæng mellem alle Ting i Verden, kan man
drage Slutninger fra den ene til den anden. Dette kan ikke skyldes
Tingenes blotte Koexistents, men maa bero paa Fællesskab i Oprindelse,
paa fælles Afhængighed af et og samme Væsen. I selve Naturordenen ligger
der altsaa et Bevis for en højeste Aarsag. Denne Paavisning af et fælles
Ophav for alle Ting i Verden som Forudsætning for Vexelvirkningen efter
almindelige Love, mener Kant at have været den Første til at give
(\textit{primus evidentissimis rationibus adstruxisse mihi videor}).

Vi staa her overfor en Tanke, som Kant ikke blot tillagde stor Betydning
i sine første Arbejder, men som han senere atter og atter vendte tilbage
til, og som er karakteristisk for hans Maade at behandle Problemer paa.
Hvor Kant staar over for et stort Princip (som f.\ Ex.\ her den mekaniske
Naturopfattelse), dér tager han det efter dets hele Rækkevidde og søger
enten at vise, at det ved nærmere Eftertanke fører til helt andre
Konsekventser end det fra først af syntes, eller ogsaa, at hele den
Sfære, det gælder for, ikke udtrykker den absolute Virkelighed, men kun
en Side ved den. Den første Fremgangsmaade anvender han i sin tidligere
Periode, den anden i sin senere Periode. Denne Fremgangsmaade har en vis
Storhed ved sig; den udmærker sig fordelagtigt frem for den saa hyppige
Halveren af Begreberne og for den Kortsynethed, der hellere søger sig et
stakket Skjul i tilsyneladende Undtagelser end resolut tænker Principet
igennem. Men den har ogsaa sine Farer, som Kants eget Exempel senere vil
vise os.

4.\ Syv Aar efter Habilitationsskriftet og den kosmogoniske Hypoteses
Udgivelse optog Kant den samme Tankegang, der havde besjælet dem, i sit
Skrift «Einzig möglicher Beweisgrund zu einer Demonstration des Daseins
Gottes» (1762). Her udføres mere i det Enkelte den Kritik af
Fysikoteologien, som allerede var indledet i «Allgemeine Naturgeschichte
und Theorie des Himmels». Kant bebrejder den, at den anser Harmonien og
Hensigtsmæssigheden i Naturen for tilfældige og derfor søger en Aarsag
til dem uden for Naturen. Den bliver derved uvidenskabelig og har ofte
hindret Erkendelsens Fremskridt, idet Nytte og Fuldkommenhed bleve
betragtede som Tegn paa, at den Ting, der besad dem, ikke havde nogen
naturlig Oprindelse, hvorfor man undlod al videre Forsken. Og selv om man
gaar ind paa den fysikoteologiske Tankegang, fører denne dog ikke til
hvad den religiøse Tro fordrer. Den forklarer nemlig kun Elementernes
Ordning, ikke deres Oprindelse. Den Guddom, hvis Tilværelse kan sluttes
efter denne Tankegang, bliver da kun Verdens Bygmester, ikke Verdens
absolute Ophav\footnote{Det har sin Interesse at fastholde, at Kants
Kritik af Fysikoteologien allerede stammer fra 1755 og 1762. Ti den
stemmer i Hovedtrækket med Humes Kritik i «Dialogues on natural
religion», der udkom 1779, to Aar før «Kritik der reinen Vernunft»
(skønt de vare udarbejdede mange Aar tidligere). Kants Selvstændighed
staar altsaa fast. Da han senere lærte «Dialogerne» at kende, satte han
stor Pris paa dem.}.

Heller ikke de andre gængse Beviser for Guds Tilværelse finde Naade for
Kants Øjne. Allerede i Habilitationsskriftet havde han sluttet sig til
Crusius' Kritik af det ontologiske Bevis. Crusius havde
vist\footnote{Chr.\ Aug.\ Crusius: \textit{Entwurf der nothwendigen
Vernunft-Wahrheiten.} Leipzig 1745.\ §\ 235.\ (Her citeret efter 3.\
Udgave 1766.)}, at dette Bevis, paa hvilket de foregaaende metafysiske
Systemer (især Descartes' og Wolfs) havde lagt saa megen Vægt, beroede
paa en Fejlslutning. Kant fører Crusius' kritiske Tankegang videre paa
lignende Maade som senere i den bekendte Udvikling i «Kritik der reinen
Vernunft» af Forskellen mellem Begreb og Existents. --- Heller ikke det
kosmologiske Bevis lader han staa. En Slutning ud fra den i Erfaringen
givne Verden kan ikke føre til Antagelsen af et absolut nødvendigt Væsen.

Den Kritik af den naturlige Teologi, som senere i «Kritik der reinen
Vernunft» vakte saa stor Opsigt, findes altsaa allerede i «Beweisgrund».
Naar man vilde holde sig til Kants Hovedværker, vilde man overse den
betydningsfulde Kendsgerning, at adskillige af de vigtigste Afsnit ere
blevne til langt tidligere.

Dog havde Kant i 1762 ikke opgivet ethvert Forsøg paa teoretisk
Begrundelse af den naturlige Teologi. Han anser vel ikke Beviser for
nødvendige, da den religiøse Tro har andet Grundlag end Bevisførelse. Men
da han anser det for værdifuldt at prøve, hvor langt man kan komme ad
Fornuftens Vej, opstiller han følgende Ændring af det ontologiske Bevis:
«Al Mulighed forudsætter noget Virkeligt, hvori og hvorved alt Tænkeligt
er givet. Der gives da en Virkelighed, ved hvis Ophævelse al indre
Mulighed i det Hele vilde falde bort. Men det, hvis Ophævelse eller
Benegtelse tilintetgør al Mulighed, er absolut nødvendigt. Altsaa er der
Noget, som existerer med absolut Nødvendighed.»\footnote{\textit{Einzig
möglicher Beweisgrund zu einer Demonstration des Daseyns Gottes.}
Königsberg 1763 [o:\ 1762].\ p.\ 29. --- Allerede i Habilitationsskriftet
findes denne Tankegang (Sectio II.\ Prop.\ 7), kun i mere dogmatisk Form.}
Meget tydelig er denne Tankegang ikke, og Kant giver den heller ikke som
Bevisførelse, men kun som en Bevisgrund, den eneste, han anser for mulig.
Kan den ikke føre til Maalet, siger han, saa kan ingen. --- At den ikke
fører til Maalet, ligger i, at selv om enhver Mulighed maa have sin Grund
i et Virkeligt, behøver dette Virkelige dog ikke at være absolut
nødvendigt; det kan selv have sin Grund i et andet Virkeligt --- og
saaledes fremdeles.

Hvad der hovedsageligt interesserer Kant ved denne forsøgte Omstøbning af
det ontologiske Bevis, er dog, at den vilde gøre en fuldstændig Afslutning
mulig af den fra hans tidligere Skrifter kendte Tankegang, der fra den
mekaniske Natursammenhæng førte til en Enhedsgrund for Alt. Hvis hver af
Tingene i Verden havde sin særskilte Nødvendighed, vilde det være et
Tilfælde, om de passede saaledes sammen, at der kunde komme et Hele ud af
dem. For at Vexelvirkning skal være mulig, maa Tingene fra først af
udgøre et System og alle være afhængige af en fælles Grund. Der maa være
et Princip, i hvilket alt Forskelligartet er forbundet, og i hvilket alt
det Mangfoldige er i Enhed. Denne Enhedsgrund, af hvilken Tingenes første
Muligheder udspringe, maa tillige være Grunden til al Visdom og Godhed:
kun derved bliver det muligt, at Tingenes lovmæssige Samvirken kan stemme
med, hvad Visdommen og Godheden fordre.

Dette er det dristigste Sving, Kants Tænkning har taget i spekulativ
Retning. Langt ud over det sædvanlige antropomorfistiske Gudsbegreb søger
han tilbage til et Grundlag, af hvilket baade den Visdom, der kundgør sig
i Verden, og de Elementer, der samvirke i Verden, udspringe. I sin Iver
efter at drage den sidste Konsekvents af det Grundfaktum, fra hvilket han
gaar ud, --- nemlig den mekaniske Natursammenhæng og den indre
Sammenhæng, den forudsætter mellem alle Verdenselementer, --- taber hans
Tanke sig tilsidst i et mystisk Dyb. Var han vedbleven at færdes dernede
i Stedet for at vende tilbage til den rationelle Bevidstheds Dagslys,
vilde han i Skumringen blandt de faa Tænkere, som færdedes der, have faaet
Øje paa Jakob Bøhme og Spinoza. Ti naar Kant første Gang han indførte den
Tankegang, hvis sidste Konsekvents han her drog, med Selvfølelse havde
erklæret, at han var dens første Ophavsmand, saa var det ikke historisk
rigtigt. I mystisk Form findes den hos Bøhme, og i klar Forbindelse med
den mekaniske Naturopfattelse og som dennes Konsekvents findes den hos
Spinoza, hvis hele System egentlig er en Gennemførelse af den Tanke, at
Aarsagssammenhængen er uforstaaelig, naar de enkelte Ting i Verden
betragtedes som selvstændige og uafhængige Væsener (Substantser i
strengeste Betydning). I sin Afhandling «de emendatione intellectus»
opstiller Spinoza netop den lovmæssige Sammenhæng mellem Fænomenerne som
det Grundfaktum, fra hvilket Tænkningen skal tage sit Udgangspunkt.
Spinozas ene Substants svarer til Kants Enhedsgrund. --- Kant har aldrig
kendt Spinoza paa første Haand og vilde sikkert være bleven højligt
forbavset over den Overensstemmelse, han var kommen i med den ensomme
Tænker, for hvem Forstaaelsens Tid endnu ikke var begyndt, da Kant skrev
sit Habilitationsskrift og sin «Beweisgrund».

I vore Dage have to aandfulde Tænkere under Bestræbelsen for at hævde en
idealistisk Verdensanskuelse uden at slaa af paa den mekaniske
Naturopfattelses Konsekvents, udviklet en lignende Tankegang som den,
hvis første Ophavsmand Kant mente at være. I Fechner's og Lotze's
filosofiske Anskuelser spiller den en vigtig Rolle. Især Lotze har paa
skarpsindig Maade grundet sin idealistiske Monisme paa Analyse af
Mekanismens Begreb\footnote{\textit{Drei Bücher der Metaphysik.} Leipzig
1879.\ p.\ 135 ff.\ (Allerede i \textit{Mikrokosmus.} Drittes Buch.\ Kap.\
5.) I den nyeste Tid har Fr.\ Paulsen sluttet sig til denne
Betragtningsmaade (\textit{Einleitung in die Philosophie.} Berlin 1892.\
p.\ 212--224).}.

5.\ En Tanke af saadan Dybde og Interesse kunde man ikke tænke sig at
Kant igen kunde lade falde, naar han én Gang havde undfanget den. Den
drager sig da ogsaa under forskellige Former hen gennem hans senere
Filosofi, skønt det forandrede Synspunkt giver den en anden Stilling og
Form. --- Samme Aar, som «Beweisgrund» blev til, indtraadte der (som det
senere skal vises) en vigtig Ændring i Kants Interesse. Denne vendte sig
fra de objektive Problemer til Undersøgelse af Metoden for deres
Behandling. Det var ad denne Vej, Kant nærmede sig til det tidligere
omtalte Vendepunkt, i hvilket den kritiske Filosofi i Principet blev til.
I selve det Skrift, i hvilket dette Princip udtaltes (\textit{Dissertatio
de mundi sensibilis et intelligibilis forma atque principiis.} 1770) møde
vi (i Kap.\ 4) den fra de ældre Skrifter bekendte Tankegang igen.
Foruden den Sammenhæng imellem Tingene i Verden, som fremtræder i Tidens
og Rummets Former, postuleres der et objektivt Princip for den reale
Vexelvirkning, hvorved en Verdenshelhed bliver mulig. Spørgsmaalet lyder:
«Hvorledes er det muligt, at Substantser kunne staa i
Vexelvirkningsforhold til hverandre og derved høre til en og samme
Helhed, som vi kalde Verden?» (§\ 16). Og der svares, at det kun er
muligt, naar de alle have ét Udspring og altsaa ikke ere Substantser i
absolut Forstand\footnote{\textit{Totum e substantiis necessariis est
impossibile} (§\ 18). --- \textit{Unitas in conjunctione substantiarum
universi est consectarium dependentiæ omnium ab Uno} (§\ 20).}.

Da Kant oversendte Dissertationen til Lambert, udtalte han (Brev til
Lambert af 2.\ Septbr.\ 1770), at det Kapitel, i hvilket denne Tankegang
findes, kunde forbigaas som uvæsentligt (\textit{unerheblich}). Sagen
var, at et helt andet Problem nu havde skudt sig i Forgrunden, idet Kant
var kommen til den Overbevisning, at Rum og Tid kun ere subjektive
Anskuelsesformer, saa at Tingene, for saa vidt vi opfatte dem i Rum og
Tid, kun kendes som Fænomener. Endnu antog han vel, at vi ved
Aarsagsbegrebet og andre Forstandsbegreber kunne erkende Tingene i sig
selv (Noumener). Men strax efter Dissertationens Fremkomst følte han dog
store Vanskeligheder ved denne Antagelse. Ved at forfølge dem kom han da
(som vi senere noget nøjere skulle vise) til det Resultat, at al vor
videnskabelige Erkendelse kun angaar Fænomener, ikke Tingene i sig selv.

Det er klart, at han nu ikke længere paa samme Maade som før kunde slutte
fra Naturmekanismen til en objektiv Enhedsgrund. Principet for den faste
Sammenhæng i Fænomenernes Verden blev et rent subjektivt Princip,
Bevidsthedens forbindende Kraft. Verdensanskuelsens Enhedspunkt laa nu
ikke længere i en mystisk Urgrund, men i selve det erkendende Subjekt.
Det hedder saaledes i en Optegnelse, der synes at stamme fra 70'erne
(Mellemrummet mellem Dissertationen og «Kritik der reinen Vernunft»):
«Expositionen af det, som er givet, beror (naar man betragter Stoffet som
ubestemt) paa Grunden til al Relation og til Sammenkædningen af
Forestillingerne (eller Fornemmelserne). Sammenkædningen grunder sig
\ldots\ ikke paa det blotte Fænomen, men er en Forestilling om den
Bevidsthedens indre Handling: at forbinde Forestillinger, ikke blot at
stille dem jevnsides hverandre for Anskuelsen, men gøre et Hele ud af dem
hvad Indhold angaar \ldots\ Expositionen af Fænomenerne er altsaa
Bestemmelsen af den Grund, hvorpaa Fornemmelsernes Sammenhæng i dem
beror.»\footnote{\textit{Lose Blätter aus Kants Nachlass.} Mitgetheilt
von Rudolf Reicke.\ Erstes Heft.\ Königsberg 1889\ p.\ 16.}

Med denne Tanke var den kritiske Filosofi fuldendt. Den hviler paa den
Grundtanke, at hele vor Verdensopfattelse bærer Præget af vor Aands
Virkemaade. Vor Fornufts Love gælde for alle mulige Fænomener, fordi
disse kun erkendes af os, derved at vor Fornuft paa sin Vis og efter sine
Love anskuer og tænker dem. Men Kontinuiteten i Kants Tænken viser sig nu
her, ti ogsaa for den kritiske Filosofi er Aarsagssammenhængen (foruden
Sammenhængen i Rum og Tid) det Grundfaktum, fra hvilket Slutningen til
Enhedsprincipet drages, --- kun at Enhedsprincipet, som sagt, nu er
subjektivt, ikke objektivt. Grundbegrebet i «Kritik der reinen Vernunft»,
Syntesen som Udtryk for Erkendelsesvirksomheden paa alle dens Trin, er
det subjektive Genbillede til den objektive Enhedsgrund, Kant i sine
Ungdomsarbejder søgte tilbage til. Fra at søge den Grund, der holder
Verden sammen og gør den til et objektivt Hele, er han nu gaaet over til
at søge den Grund, der holder Verdensbilledet sammen og gør det til et
subjektivt Hele. I sit inderste Væsen som forbindende Enhedsvirken, som
Syntese er den erkendende Aand da et Forbillede paa Alt, hvad der skal
kunne erkendes, da dette stedse maa have et Enhedspræg. Som Kant siger:
«Ich bin das Original aller Objecte», eller: «Das Gemüth [o:\
Bevidstheden] ist sich selbst das Urbild der
Synthesis»\footnote{\textit{Lose Blätter} I.\ p.\ 19.\ 20.}. --- Ikke
længere Substantsen, men Syntesen er Grundbegrebet. Derved betegnes den
kritiske Filosofis Modsætning til den dogmatiske.

6.\ I «Kritik der reinen Vernunft» findes der, hvad Kritiken af de saa
kaldte Beviser paa Guds Tilværelse angaar, intet egentligt Nyt i Forhold
til de førkritiske Skrifter. Det Nye er her et Negativt: Udeladelsen af
den Tankegang, som i sin Tid (endnu i Dissertationen) for Kant havde
dannet en Bro mellem den videnskabelige Erkendelse og den religiøse Tro.
Denne Bro faldt bort, da Forstandsbegrebernes Anvendelse indskrænkedes
til Erfaringens Verden. Kun ved et Spring er det nu muligt at komme fra
Erfaringens betingede Verden til det Ubetingede, den religiøse Tro søger,
og dette Spring sker i Kraft af praktiske, moralske Motiver, ikke af
Erkendelsestrang. Erfaringen, som altid er betinget, kan ikke begrunde
Tanken om en ubetinget Aarsag\footnote{\textit{Kritik der reinen
Vernunft} 2.\ Aufl.\ p.\ 665. --- Den Tankegang, ved hvilken Kant i sin
Tid (se §\ 4) havde forsøgt at omstøbe det ontologiske Bevis, har han nu
ogsaa opgivet som bevisende. Men Ideen om det absolute Væsen som alle
Muligheders Grund er dermed ikke falden bort. Den fremtræder som «det
transscendentale ideal», den højeste Maalestok for Alt, hvad der tillægges
Realitet. Se \textit{Kritik der reinen Vernunft.} 2.\ Aufl.\ p.\ 600;
606. Allerede i «Dissertationen» er den højeste teoretiske Ide
\textit{exemplar aliquod, omnium aliorum, quoad realitates, mensura
communis} (§\ 9). Her er altsaa foregaaet en med «Enhedsgrunden»s
Subjektivering analog Metamorfose.}. Hermed er Slutningen fra
Aarsagssammenhængen i Verden til en absolut Enhedsgrund
bortfalden\footnote{Vexelvirkningsproblemet omtales dog \textit{Prolegomena}
p.\ 98 f.\ og \textit{Kritik der reinen Vernunft.} 2.\ Aufl.\ p.\ 265
Note, 292 f., 428, og paa en Maade, der viser, at Kant ikke har glemt sin
gamle Tanke. Smign.\ ogsaa \textit{Reflexionen Kants.} II.\ p.\ 219 f.\
\textit{Lose Blätter} I.\ p.\ 274.}.

Men den store Interesse for Aarsagserkendelsen og den konsekvente
Gennemførelse af den i dens strengt videnskabelige Form falder ikke bort
med Indskrænkningen af dens Anvendelighed til Fænomenernes Verden. Den
genfindes i den Maade, hvorpaa Kant fører alle de forskellige
Grundsætninger, der indeholde Betingelserne for en mulig Erfaring,
tilbage til Kontinuitetens Lov. De fire Klasser af Grundsætninger, der i
Kants System svare til de fire Klasser af Kategorier (lige som disse igen
svare til de fire Klasser af logiske Domme), resumeres i fire Motto'er:
1) \textit{In Mundo non datur hiatus.} 2) \textit{In Mundo non datur
saltus.} 3) \textit{In Mundo non datur casus.} 4) \textit{In Mundo non
datur fatum.} Og som fælles Karaktermærke for dem alle gælder det, at de
udelukke, hvad der vilde afbryde den kontinuerlige
Sammenhæng\footnote{\textit{Kritik der reinen Vernunft,} 2.\ Aufl.\ p.\
282.}. Kant har paa dette mærkelige Sted samlet hele sin Erkendelsesteori
i Kontinuitetens Lov. Havde han lagt dette Synspunkt til Grund i Stedet
for at følge den logiske Systematik, han til sin Glæde havde faaet
konstrueret, vilde hans Grundtanker være komne mere til deres Ret. Der er
Spor af, at han har tænkt paa at gaa denne Vej, idet han i en Optegnelse
fra 70'erne prøver at udlede Aarsagsloven af den almindelige
Kontinuitetslov: «At alt Tilfældigt, eller hvad der opstaar, maa have sin
Grund, kommer af, at der uden Princip ikke vilde være nogen Kontinuitet i
Fænomenerne, og uden Regel ingen Identitet.»\footnote{\textit{Reflexionen
Kants.} I.\ p.\ 308 (Nr.\ 1074). Smlgn.\ p.\ 514 (Nr.\ 1747--50).}

Den Tanke, som her ligger til Grund, er, at hvad vi søge i Forstaaelsen
af et Fænomen, ikke er en blot udvortes Sammenstilling af det med andre
Fænomener, men en saadan nøje og bestemt Forbindelse mellem dem, at det
Fænomen, vi søge at forstaa, kan staa som Fortsættelse af de foregaaende
Fænomener og sammen med dem danne en kontinuerlig Række. Det
Usammenhængende og Isolerede er os uforstaaeligt. Denne vigtige
Betragtning har Kant kun antydet; men den ligger som Spire i den
Opfattelse og Anvendelse af Aarsagsbegrebet, som han fra sin Ungdom af
tog Ordet for. Den skarpsindigste af hans nærmeste Efterfølgere har klart
udtalt denne Tanke\footnote{Salomon Maimón: \textit{Versuch über die
Transscendentalphilosophie.} Berlin 1790.\ p.\ 140: «Die Ursache der
Erscheinung [aufzusuchen, heisst] das Stetige in derselben aufzusuchen,
und die Lücken unserer Wahrnehmungen auszufüllen, um sie dadurch zu
Erfahrungen zu machen. Denn was versteht man sonst in der Naturlehre
unter dem Wort Ursache als die Entwickelung einer Erscheinung und
Auflösung derselben, so dass man zwischen ihr und der vorhergehenden
Erscheinung die gesuchte Stetigkeit finde.»}.

Den ubrudte Natursammenhæng hævder Kant energisk for Fænomenalverdenens
Vedkommende. I det interessante Afsnit «Die Disciplin der reinen Vernunft
in Ansehung der Hypothesen» fordrer han, at ingen Forklaringsgrunde
anvendes, som ikke efter bekendte Love hænge sammen med de givne
Fænomener. Det er Principet om \textit{vera causa}, som han selv paa en
saa storartet Maade havde anvendt i sin kosmogoniske Hypotese. Det har
vedblivende sin Plads i hans Metodelære. Og han fordrer det særligt
anvendt, hvor der har vist sig størst Fristelse til at krænke det, nemlig
ved Forklaringen af Ordenen og Hensigtsmæssigheden i Naturen. «Naturens
Orden og Hensigtsmæssighed», siger han, «maa stedse forklares af
Naturgrunde og efter Naturlove, og her ere selv de vildeste Hypoteser,
naar de kun ere naturlige, taaleligere end overnaturlige, d.\ v.\ s.\ end
Beraabelsen paa en overnaturlig Ophavsmand, som man antager til dette
Behov»\footnote{\textit{Kritik der reinen Vernunft.} 2.\ Aufl.\ p.\ 800 f.}.

7.\ Endnu mere afgjort vender Kant i sit sidste Værk paa den teoretiske
Filosofis Omraade tilbage til sin Ungdoms store Tanke. «Kritik der
Urtheilskraft» (1790) stiller sig paa en Maade samme Opgave som
«Allgemeine Naturgeschichte» og «Beweisgrund», nemlig Undersøgelsen af,
hvorledes Skønheden og Hensigtsmæssigheden i Verden er at forklare. Kant
kommer efter mange kritiske Betænkeligheder til det Resultat, at den
Forskel, vi i vor Betragtningsmaade af Naturen gøre mellem mekanisk og
teleologisk Opfattelse, ikke berettiger til den Antagelse, at en analog
Principernes Dobbelthed skulde bestaa i Tingenes Væsen. Han søger at vise,
at naar vi over for visse Naturfænomener paa Grund af deres specielle
Ejendommelighed ledes til at antage et Formaal som forklarende Grund,
medens vi ved simplere Fænomener ikke føle denne Trang, --- saa har dette
sin Grund i vor Forstands Natur, og vi tør ikke tillægge det absolut
Betydning. Det saakaldte teleologiske Problem er for Kant kun en Form for
det almindelige Problem, som al vor Erfaring stiller os: hvorledes det
Mangfoldige kan være forbundet til Enhed gennem fælles Love, --- det
Problem, der har fulgt Kant gennem hele hans filosofiske Udviklingsgang.
Og skønt han --- i det Mindste hvad Udtryksmaade angaar --- gør den
naturlige Teologi saa store Indrømmelser som muligt, peger han dog som
sit egentlige Resultat hen til den Mulighed, at «den fysisk-mekaniske og
den teleologiske Forbindelse i Naturens os ubekendte Grund kunne hænge
sammen i ét Princip, kun at vor Fornuft ikke er i Stand til at forene dem
i et saadant Princip»\footnote{\textit{Kritik der Urtheilskraft.} §\ 70
(kfr.\ §\ 76, 80, 88).}. Den fra Ungdomsskrifterne kendte Tanke om et
fælles Grundlag for Tingenes sidste Muligheder og for den guddommelige
Visdom viser sig her ingenlunde at være opgiven, skønt den træder frem
med endnu større Forsigtighed end før. Kant har beskrevet en Bue fra
sine første naturfilosofiske Afhandlinger gennem den kritiske Skærsild
til de dybsindige Antydninger, med hvilke hans Tankearbejde ender. Trods
de ændrede Synspunkter bevarer der sig dog en Kontinuitet i
Verdensopfattelsen, som det her har sin særegne Interesse at lægge Mærke
til.

\section*{II.\\ Analyse og Konstruktion.}

\noindent 8.\ Vi vende nu tilbage til et tidligere Punkt i Kants
Udvikling for at forfølge en anden Tankerække hos ham, lige som vi have
søgt at følge Aarsagsbegrebet. At de forskellige Tankerækker gribe ind i
hverandre under Udviklingen, følger af sig selv; men derfor kan det
alligevel have sin Fordel at udrede dem hver for sig. Kontinuiteten
træder derved klarere frem. Og paa de vigtigste Punkter ville vi ikke
undlade at undersøge deres Vexelvirkning.

Aaret 1762 (og Begyndelsen af det følgende Aar) er et af de frugtbareste
i Kants Forfatterliv. Ikke mindre end fire betydningsfulde Afhandlinger
skrive sig fra dette Aar, som allerede ved denne stærke Produktion drager
Opmærksomheden til sig. Det ene af disse Arbejder kende vi allerede:
«Einzig möglicher Beweisgrund». Det behandlede Aarsagsbegrebet i dets
Betydning for Naturvidenskaben og som forbindende Led mellem
Naturerkendelse og Gudserkendelse. Det indeholdt som særligt Exempel en
kort Oversigt over den syv Aar forud opstillede kosmogoniske Hypotese. De
andre tre Skrifter («Die falsche Spitzfindigkeit der syllogistischen
Figuren»; «Ueber die Deutlichkeit der Grundsätze der natürlichen Theologie
und Moral»; «Versuch den Begriff der negativen Grössen in die Weltweisheit
einzuführen»)\footnote{De fire Skrifter synes (smign.\ B.\ Erdmann:
\textit{Reflexionen Kants.} I.\ p.\ XVII ff.) at være udkomne i følgende
Orden: 1) Spitzfindigkeit; 2) Beweisgrund; 3) Deutlichkeit; 4) Negative
Grössen. --- For Forholdet mellem deres Tankeindhold er Affattelsens og
Udgivelsens Rækkefølge naturligvis ikke afgørende.} drøfte tilsyneladende
helt andre Spørgsmaal. Og dog bleve de Undersøgelser, der anstilledes i
dem, af stor Betydning netop for Opfattelsen af det Begreb, paa hvis
Betydning for Naturvidenskab og naturlig Teologi Kant endnu i
«Beweisgrund» havde lagt saa stor Vægt. Der har i det nævnte Aar (hvis
Produktivitet maaske netop forklares derved) bevæget sig to Tankerækker i
Kants Bevidsthed, hvis Sammenstød maatte stille ham over for et vigtigt
Problem.

I sin Kritik af Slutningsfigurerne gaar Kant ud fra det Synspunkt, at
enhver Dom er en Sammenligning. At dømme vil sige at sammenligne et
Kendemærke med en Ting og som et Resultat af denne Sammenligning enten
tillægge eller frakende Tingen
det\footnote{«Etwas als ein Merkmal mit einem Dinge vergleichen heisst
urtheilen» (\textit{Ueber die falsche Spitzfindigkeit.} §\ 1). ---
Smlgn.\ «Träume eines Geistersehers» (1766): «Unsere Vernunftregel gehet
nur auf die Vergleichung nach der Identität und dem Widerspruche». (2.\
Th.\ 3.\ Hauptst.) (Kehrbachs Udgave, p.\ 64.)}. En lignende Opfattelse
af Tænkningens Natur ligger til Grund i Afhandlingen «om Grundsætningernes
Tydelighed». Her indskærpes Forskellen mellem den filosofiske og den
matematiske Metode. Filosofiens Metode er Analyse, Matematikens
Konstruktion. Matematiken kan strax danne færdige Begreber, med hvilke
der kan opereres. Men i Filosofien, der henter sit Stof fra Erfaringen,
bliver Begrebsbestemmelsen først færdig, naar Analysen af det Givne er
fuldendt. Her er da den Vanskelighed, hvorledes man kan overbevise sig
om, at Analysen er ført tilstrækkeligt vidt, saaledes at der ikke er
flere Kendemærker at opdage. I Matematiken danne derfor Definitioner
Begyndelsen, medens de i Filosofien først kunne komme i Slutningen. Den
tidligere Filosofis Ufuldkommenheder udledes af, at man har opereret med
ufærdige Begreber og indladt sig paa forhastede Konstruktioner. Som et
Exempel paa et Begreb, man i den dogmatiske Filosofi trygt har opereret
med, som om det var færdigt og afsluttet, nævnes Begrebet Aand (tænkende
Substants), der spillede saa stor en Rolle hos Descartes, Leibniz og
Wolf. Det var vilkaarligt konstrueret, ikke grundet paa gennemført
Analyse\footnote{«Träume eines Geistersehers» er egentlig kun en nærmere
Udvikling af dette Exempel for at vise, hvad det fører til at operere med
ufærdige Begreber. Begrebet Aand (i spiritualistisk Betydning) erklæres
her for et tilsneget Begreb. Resultatet af det aandfulde Skrift er: «Die
Pneumatologie der Menschen kann ein Lehrbegriff ihrer nothwendigen
Unwissenheit in Absicht auf eine vermuthete Art Wesen genannt werden»
(1.\ Th.\ 4.\ Hauptst.). (Kehrbachs Udgave, p.\ 43.) Smlgn.\ hermed Kants
«Nachricht von der Einrichtung seiner Vorlesungen» (1765--66), hvor det
hedder, at den empiriske Psykologi slet ikke spørger, om Mennesket har en
Sjæl. --- Som man ser, var ikke blot Kritiken af Teologien, men ogsaa
Kritiken af den rationale [o:\ spiritualistiske] Psykologi væsentligt
færdig paa dette Punkt af Kants Udvikling.}. At gennemføre Analysen af
indviklede Erkendelser erklærer Kant for langt vanskeligere end at
forbinde simple Erkendelser ved Syntese og bygge Slutninger paa dem.
Metafysiken er derfor den vanskeligste af alle Videnskaber --- men der er
da heller ikke endnu skrevet nogen Metafysik! Det vil endnu vare meget
længe, inden man i Metafysiken kan gaa syntetisk frem. Det vil først
kunne ske, naar Analysen har hjulpet til tydelige og udførligt udviklede
Begreber.

En Konsekvents maatte her ligge nær: Kun saadanne Forhold ere
forstaaelige, hvor Tænkningen ad Analysens Vej kan føre fra det ene Led i
Forholdet over til det andet. Denne Konsekvents drog Kant dog ikke i
Afhandlingen «om Tydeligheden». Derimod faar det fjerde Skrift,
Afhandlingen om de negative Størrelser i Filosofien, sin store Interesse
ved, at denne Konsekvents drages med fuld Klarhed. Idet Kant i dette
Skrift indskærper den af alle Tiders spekulative Filosofer saa ofte
oversete Forskel mellem logisk Negation og real Modsætning, kommer han
naturligt ogsaa til at undersøge saadanne Tilfælde, hvor Noget ophæves,
fordi noget Andet indtræder; og her staar han da over for
Aarsagsproblemet. Ad Analysens eller Sammenligningens Vej kan man ikke
vise Nødvendigheden af Sammenhængen mellem det Enes Indtræden og det
Andets Ophør, saalidt som i det Hele af, at Noget sker, fordi noget Andet
sker. Kant lader Problemet henstaa uløst. Kun saa meget er han sikker
paa, at Modsigelsens Grundsætning, til hvilken Dogmatismen vilde føre al
Begrundelse tilbage, ingen Forklaring
giver\footnote{\textit{Versuch den Begriff der negativen Grössen in die
Weltweisheit einzuführen.} (Udtalelserne om Aarsagsproblemet findes i
Slutningsanmærkningen). --- Smlgn.\ «Träume» (2.\ Th.\ 3.\ Hauptst.):
«Unsere Vernunftregel gehet nur auf die Vergleichung nach der Identität
und dem Widerspruche. So ferne aber etwas eine Ursache ist, so wird
dadurch etwas gesetzt, und es ist also kein Zusammenhang vermöge der
Einstimmung anzutreffen; wie denn auch, wenn ich eben dasselbe nicht als
eine Ursache ansehen will, niemals ein Widerspruch entspringt».
(Kehrbachs Udgave, p.\ 64.)}.

Kant har her stillet Aarsagsproblemet, til Dels i lignende Udtryk som de,
i hvilke Hume allerede havde stillet det. Men Kant havde jo egentlig
allerede havt med Aarsagsproblemet at gøre i sine tidligere Skrifter,
senest i «Beweisgrund». Hvad der laa til Grund for den Tankegang, der for
ham forbandt Natur- og Gudserkendelsen, og som han mente først at have
ført ind paa (se §\ 3), var jo Umuligheden af at forstaa Vexelvirkningen,
Aarsagssammenhængen mellem Tingene i Verden, naar man ikke antog en
fælles Grund for dem. Saa længe de staa i deres Forskellighed, er
Aarsagsforholdet mellem dem uforstaaeligt. Det, der er sket i de
mærkelige Slutningsudtalelser i Afhandlingen om de negative Størrelser,
er da egentlig ikke en Opstilling af et helt nyt Problem, men
Omsætningen af et Problem, der hidtil behandledes i metafysisk-objektiv
Form, i \emph{erkendelsesteoretisk-subjektiv} Form. Det er forstaaeligt,
at Kant, da han nu paa ny havde gennemarbejdet de Synsmaader, han syv Aar
forud første Gang havde fremstillet, og da han saa samtidigt kom ind paa
rent logiske og metodiske Undersøgelser, ved disse to Tankerækkers
Sammenstød maatte komme til at se, at Aarsagsforholdet ikke blot stiller
et objektivt, men ogsaa et subjektivt Problem. Af dette Aars saa
intensive og produktive Tankearbejde er denne Omsætning den vigtigste
Frugt, ganske vist en Frugt, der endnu ikke kunde komme til fuld
Modenhed. Det Begreb, som han selv i sin Naturforsken havde anvendt paa
saa genial en Maade, og som endnu for nylig havde slaaet Bro for ham
mellem Religionens og Naturvidenskabens Regioner, det viste sig nu
pludseligt at indeholde et saa stort Problem, naar man opkastede det
simple Spørgsmaal, hvorledes man kommer fra det ene Led i det Forhold,
det angiver, til det andet.

At der var en saadan Sammenhæng for Kant mellem Aarsagsproblemet i dets
metafysiske og i det erkendelsesteoretiske Form, ses tydeligt af
Udtalelsen i «Träume eines Geistersehers» (i Slutningskapitlet). Her
gøres paa ny opmærksom paa Aarsagsproblemet, og det sker i Sammenhæng med
det almindelige Problem om de aandelige Væseners Natur og deres Forhold
til Materien. Det vises, hvorledes Problemerne begynde indenfor
Spekulationen, hvor man trygt opererer med «Grundforholdene» (Aarsag og
Virkning, Substants og Handling), men at man, naar man bliver ved at
filosofere, ender med at finde Vanskeligheder i selve Grundforholdene. Og
denne Sammenhæng mellem det metafysiske og det erkendelsesteoretiske
Problem træder endnu tydeligere frem i Brevet til Mendelssohn af 8.\
April 1766, idet her det almindelige Spørgsmaal, hvorledes Noget kan være
Aarsag, udløser sig af det mere specielle, hvorledes en aandelig
Substants kan have Evne til at virke og lide i Forhold til Materien.

Jeg finder i Henhold til det Foregaaende Kontinuiteten i Kants Udvikling
paa et andet Punkt end Friedrich Paulsen. Denne Forsker lægger Vægten
paa, at Kant ved nærmere Undersøgelse af det metafysiske Gudsbegreb, i
Følge hvilket Gud skulde være Indbegrebet af alle Realiteter, opdagede,
at der jo er Realiteter, som staa i reelt Modsætningsforhold til
hinanden og altsaa ikke lade sig forene. Afhandlingen om de negative
Størrelser skulde da være en Specialundersøgelse, der fremkaldtes ved de
teologiske Betragtninger i
«Beweisgrund»\footnote{\textit{Entwickelungsgeschichte der Kantischen
Erkenntnisstheorie.} p.\ 64 ff.}. At der er en Sammenhæng mellem
«Beweisgrund» og Afhandlingen om de negative Størrelser, negter jeg ikke.
Allerede i «Beweisgrund» hævdes, som Paulsen har vist, den vigtige
Forskel mellem logisk Negation og real Modsætning. Men Hovedsagen er for
mig den ejendommelige og pludselige Vending, Kant gør i
Slutningsanmærkningen til det sidstnævnte Skrift. Han kunde have drøftet
de negative Størrelsers Problem udførligt og grundigt, uden just at drage
den Konsekvents, han der drager. Det er denne Vending, jeg kun kan
forklare mig ved, at Aarsagsforholdet netop var Kant nærværende som det
vigtigste i al Erkendelse. Det maatte da komme til et Sammenstød.

9.\ Kants hele Stemning ved Udgangen af det paa Tankearbejde saa rige Aar
var afgjort antidogmatisk. Fra tryg Opereren med hidtil anerkendte
Begreber var han nu bleven ført til at undersøge nogle af de vigtigste af
disse Begreber og havde fundet dem uklare og ufærdige. Han følte nu store
Vanskeligheder ved det, som Andre --- og han selv hidtil --- fandt let.
Med Ironi vender han sig mod «de grundige Filosofer, hvis Tal stadigt
forøges», med Anmodning om at løse ham disse simple Spørgsmaal, han var
stanset ved. Endnu stærkere og kaadere kom denne Stemning frem i «Träume
eines Geistersehers», der spotter over de «Luftbygmestre», som konstruere
deres Bygninger af tilsnegne Begreber, og overfor hvem der ikke er Andet
at gøre end at vente, til de have drømt ud.

Selve den stærke Modsætning mellem Konstruktion og Analyse betegner et
Brud med den foregaaende Tids Filosofi. Kant fordrer et helt nyt Grundlag
skaffet til Veje, før Tiden til filosofisk Systematisering kunde komme.
Hans Opmærksomhed var fra nu af rettet mod Metoden, mod Undersøgelsen af
Erkendelsens Grænser. Han er sikkert\footnote{Som Benno Erdmann har
paavist i Indledningen til hans Udgave af \textit{Prolegomena,} p.\ LXXVI
f.\ og udførligere i \textit{Reflexionen Kants.} II.\ p.\
XXXV--XLVII.}\ allerede paa denne Tid slaaet ind paa den antinomiske
Metode for at paavise Grænsen for Erkendelsen, idet han fandt en saadan
Grænse, hvor der angaaende samme Problem kunde begrundes selvmodsigende
Sætninger. «Jeg forsøgte», siger han senere\footnote{\textit{Reflexionen
Kants.} II.\ p.\ 4 (Nr.\ 4). Hermed stemmer en Udtalelse i Brevet til
Garve af 21.\ September 1798, at det var den rene Fornufts kosmologiske
Antinomier, der «først vakte ham af den dogmatiske Slummer og drev til
Kritik af Fornuften for at hæve Skandalen af Fornuftens tilsyneladende
Modsigelse med sig selv». (Brevet er trykt hos Alb.\ Stern: \textit{Ueber
die Beziehungen Garve's zu Kant.})}, «for fuldt Alvor at bevise Sætninger
og det Modsatte af dem, ikke for at hævde Skepticisme, men fordi jeg
formodede, at der var en Forstandsillusion til Stede, og for at opdage,
hvori den bestod». I et Brev til Lambert (af 31.\ December 1765) siger
han, at hans Bestræbelser hovedsageligt gaa ud paa Metafysikens
ejendommelige Metode. Han har nu efter flere Aars Overvejelse fundet en
Metode, som han anvender for at undgaa indbildt Viden. Ved enhver
Undersøgelse ser han efter, hvad han maa vide for at kunne løse
Problemet, og hvilken Grad af Erkendelse det Givne tilsteder; herved
bliver hans Dom vel ofte mere begrænset, men ogsaa sikrere end det
sædvanligvis er Tilfældet. I Overensstemmelse hermed udtaler han i
«Träume», at Metafysik er Læren om Erkendelsens Grænser, og i «Nachricht
von der Einrichtung seiner Vorlesungen», at «Critik der Vernunft» er en
Hovedopgave for ham ved hans Bearbejdelse af Logiken. Omtrent til samme
Tid hører sikkert ogsaa et Fragment, i hvilket det hedder: «Die
Metaphysik ist die Kritik der menschlichen Vernunft \ldots\ [sie] ist
subjectiv und problematisch»\footnote{\textit{Reflexionen Kants.} II.\
p.\ 158 (Nr.\ 507).}.

10.\ De Kantforskere have sikkert Ret, som have søgt at vise, at den
Udvikling, der foregaar i Kants Tankegang i 1762 og de følgende Aar, ikke
med Nødvendighed forudsætter en Paavirkning ude fra, men i og for sig
meget vel kan forstaas af hans forudgaaende Udvikling. Men paa den anden
Side kender jeg intet Tidsrum i Kants Udvikling, paa hvilket Udtrykket
«Vækkelse af dogmatisk Slummer» vilde passe saa godt som her. Selv Benno
Erdmann, der henlægger Vækkelsen langt senere (saa langt senere, at jeg
synes, han kommer i Konflikt med Kants Udsagn om, at Vækkelsen skete «vor
vielen Jahren»), erklærer Tiden i 60'erne for den Periode, i hvilken «de
Kantiske Tanker vare mest i Strømning». At vækkes af dogmatisk Slummer
vil netop sige, at der kommer Strømning i de hidtil faste Tanker. Ganske
vist, hvis man ved Udtrykket «Vækkelse af dogmatisk Slummer» vil forstaa
den fuldstændige Overgang til den kritiske Filosofi, saa passer det ikke
paa dette Tidspunkt. I denne strenge Betydning tager B.\ Erdmann det,
naar han først finder det anvendeligt, hvor Haabet om en
Forstandserkendelse af Tingene i sig selv «ikke blot var udryddet til
sidste Trevl, men endog blev erstattet ved en kontrært modsat
Opfattelse»\footnote{\textit{Einleitung zu Prolegomena,} p.\ XCII.}.
Dette er dog vist at forlange for meget af en Vækkelse. Kant selv har i
et Udkast til Fortale til «Kritik der reinen Vernunft» udtalt: «Det
varede længe, inden jeg fandt \emph{hele} den dogmatiske Teori dialektisk.
Men jeg søgte efter noget Vist, om ikke i Henseende til Genstanden, saa
dog i Henseende til Erkendelsens Natur og
Grænser.»\footnote{\textit{Reflexionen Kants.} II.\ p.\ 4 (Nr.\ 3). ---
Jeg lægger Eftertrykket paa Ordet «hele».} Den, der søger saaledes, er
vaagen. Vi have ogsaa hørt, at han selv paa denne Tid betragtede
Dogmatikerne som drømmende. Jeg kan ikke se Andet, end at den skarpe
Modsætning mellem Konstruktion og Analyse, den store Vægt paa Metoden og
paa Erkendelsens Grænser samt den ringe Mening om hvad der hidtil var
opnaaet i Filosofien, afgjort stiller Kant uden for Dogmatikernes Kreds
allerede paa denne Tid. Kun fra et (i Kantisk Forstand) «hyperkritisk»
Standpunkt vilde han i denne Periode kunne kaldes hvilende i dogmatisk
Slummer. Skulde Kant alligevel i Aaret 1783 have regnet Perioden i
60'erne til Dogmatismen, saa har han gjort sig selv højlig Uret. Hans
egen Definition af Begrebet Dogmatisme var jo, at den bestod i «den
Anmasselse at ville slaa sig igennem blot med en ren Erkendelse af
Begreber, efter saadanne Principer, som Fornuften forlængst har havt i
Brug, uden at undersøge den Maade, hvorpaa, og den Ret, hvormed den er
naat til dem»\footnote{\textit{Kritik der reinen Vernunft.} 2.\ Aufl.\
Vorrede, p.\ XXXV.}, eller, som det et andet Sted hedder: «den almindelige
Tillid til Metafysikens Principer uden forudgaaende Kritik af selve
Fornuftevnen»\footnote{\textit{Ueber eine Entdeckung, nach der alle neue
Critik der reinen Vernunft durch eine ältere entbehrlich gemacht werden
soll.} Königsberg 1790.\ p.\ 78.}. Efter denne Definition kunde Kant i
Tiden fra 1762 ikke kaldes Dogmatiker. Han havde ganske vist ikke, end
ikke i «Träume eines Geistersehers», opgivet Haabet om, at ihærdigt
Arbejde ad Analysens Vej en Gang kunde gøre en konstruktiv Erkendelse
mulig. Men han havde et klart Blik for Vanskelighederne og saare liden
Tillid til de gængse Principer, samt en bestemt Overbevisning om, at
Resultatet, om der kom noget, vilde være forbundet med en Begrænsning af
vor Erkendelse. Hvis han selv senere har regnet denne Periode med til den
dogmatiske Slummerperiode, saa maa «Kritiken»s Resultater have blændet
hans Blik for hans egen Fortid. --- Spørgsmaalet om Vækkelsens Tidspunkt
er det dog ikke muligt at afgøre med Sikkerhed efter hvad der foreligger.
Jeg kan kun ikke se Andet, end at Henlæggelsen til Aaret
1762--63\footnote{Til dette Tidspunkt henføres «Vækkelsen» af Kuno
Fischer, Riehl og Vaihinger (hvilken Sidste dog ogsaa antager en senere
Vækkelse). Min Motivering er dog forskellig fra disse Forskeres.} er mest
sandsynlig, skønt saa fremragende Kantforskere som Fr.\ Paulsen og Benno
Erdmann ere komne til andre Resultater. Som gentagne Gange bemærket, er
Usikkerheden om dette Spørgsmaal et Vidnesbyrd om Kontinuiteten i Kants
Udvikling.

11.\ Benno Erdmann\footnote{«Kant und Hume um 1762». \textit{Archiv für
die Geschichte der Philosophie.} I.\ p.\ 62--77; 216--230.} mener af
Herders Udtalelser fra de Aar, i hvilke han var Kants Tilhører
(1762--64), at kunne slutte, at Kant ikke paa den Tid kan være bleven
vakt af Hume. Man maatte ellers, mener Erdmann, i den Maade, hvorpaa
Herder omtaler Hume, have mærket Noget af den Begejstring, Kant utvivlsomt
maatte have indgydt sine Tilhørere for den engelske Tænker.

Denne Antagelse forekommer mig ikke tvingende, især naar man fastholder,
at Vækkelsen maa have havt en højst indirekte Karakter, saa at den
væsentligt har bestaaet i en Udløsning og fuld Udklaring af Tanker og
Tvivl, der allerede vare begyndt at arbejde sig frem. Det var, efter
Kants Udsagn, «Erindringen om Hume», der vakte ham. Altsaa Kant har
tidligere læst Hume (det er dennes «Inquiry» i tysk Oversættelse, det her
drejer sig om), men Tanken om Humes Problemstilling har først senere
paavirket ham. Det er let forstaaeligt, at denne Erindring om en Lecture,
der i sin Tid ikke havde gjort noget stærkere Indtryk, kunde komme frem
hos Kant netop i det Aar, da de to Tankerækker, der hidtil havde gaaet
Side om Side i hans Bevidsthed, --- Aarsagsforholdet som Udtryk for
forskellige Elementers Forbindelse og Tænkningen som Analyse, opererende
med Modsigelsens Grundsætning alene, --- begge optoges til ny,
indtrængende Bearbejdelse og let maatte komme til at støde sammen. Man
kan antage, at «Erindringen om Hume» enten har begunstiget dette
Sammenstød eller forstærket dets Virkninger. Men Kant har i hvert Tilfælde
neppe lige strax kunnet have fuld Bevidsthed om Betydningen af hvad der
gik for sig hos ham; en saadan kunde først komme, naar den
Udviklingsproces, der indlededes ved Vækkelsen, laa helt færdig for det
tilbageskuende Blik. Derved forstaas ogsaa, at han ikke omtaler Vækkelsen
ved Hume i sine Skrifter før længe efter og, som Erdmann selv har oplyst,
paa speciel Foranledning.

Der er ingen Grund til at antage, at Kant paa sine Forelæsninger i den
Tid, Herder hørte ham, skulde have omtalt Hume (som teoretisk Filosof)
anderledes, end han senere stedse omtalte ham, nemlig som Skeptikeren,
der danner en sund Modvægt og et gavnligt Ferment over for Dogmatismen.
Endnu manglede Kant jo oven i Købet den «højere Enhed», Kriticismen, som
han senere plejede at lade triumfere over Dogmatisme og Skepticisme, og
han maatte da overvejende betragte Hume som en Modstander, der skulde
overvindes. I de Ytringer, som Erdmann anfører af Herder, omtaler denne
Hume netop paa den Maade, vi maatte vente, nemlig som et fint Hoved, men
tillige som Erketvivler; han beskylder ham endog for «ein erstrebtes
Zweifeln». Jeg forstaar ikke, hvorledes Erdmann efter disse Ytringer kan
slutte\footnote{\textit{Archiv.} I.\ p.\ 216.}, at Kant --- naar man
støtter sig paa Herders Vidnesbyrd --- i disse Aar ikke i Hume
paaskønnede den metafysiske Tvivler, men kun Moralisten. Tvivleren i Hume
synes Herder dog at være bleven gjort tilstrækkeligt opmærksom paa!
Herder har sikkert endog overdrevet Sagen. Udtrykket «erstrebt» maa
sandsynligvis staa for hans Regning, ikke for
Kants\footnote{Et Brev fra Kant til Herder af 9.\ Maj 1767 (som jeg kun
ved en venlig Foræring fra Hr.\ Bibliotekar Dr.\ R.\ Reicke i Königsberg
blev opmærksom paa) viser, at Kant dog paa denne Tid havde en mere positiv
Opfattelse af Hume end man af Herders Ytringer skulde antage. Kant ønsker
sin unge Ven, at han maa faa den (ingenlunde følelsesløse) Ro i Stedet,
som er ejendommelig for Filosofen i Modsætning til Mystikeren, og tilføjer
derpaa: «Ich hoffe diese Epoche Ihres Genies aus demjenigen was ich von
Ihnen kenne mit Zuversicht, eine Gemüthsverfassung, die dem so sie
besitzt und der Welt unter allen am nützlichsten ist, worin Montaigne den
untersten und Hume so viel ich weiss den obersten Platz einnehmen». Denne
i Kants Mund store Ros maa efter hele Sammenhængen ikke blot gælde
«Moralisten». (Brevet er trykt i \textit{Altpreuss.\ Monatsschrift.} Bd.\
28.\ Heft 3 u.\ 4.\ 1891: \textit{Zu Herders Briefwechsel. Von Victor
Diederichs.})}. Det var efter Herders hele Naturel og aandelige Retning
ikke underligt, at Humes Tvivl kunde synes ham overdreven og vilkaarlig.
Herder havde ikke den Brug for den som Kant, der fik sine Tanker stærkere
paa Gled ved den, ja, Herder kunde vel ikke en Gang forstaa, hvilken Brug
Kant kunde have for den; hans senere Forhold til Kant lader formode det.
Hvad Herder paaskønnede i saa høj Grad hos Kant paa denne Tid, var (efter
den bekendte Udtalelse i «Humanitetsbrevene») den aabne Sans for alt
Menneskeligt, saa vel som for Naturen, den livlige Aand og den klare
Forstand. Den søgende, ad ensomme og tornede Veje vandrende Tænker havde
Herder ikke Sans for.

12.\ Ved den klare Distinktion mellem Konstruktion og Analyse havde Kant
gjort et stort Fremskridt. De dogmatiske Systemers Fejl var, at de gik
for rask til Konstruktionen, selv om de naturligvis ikke ganske kunde
undvære Analysen, der --- om end nok saa ufuldstændig og ufærdig --- i
Virkeligheden altid danner Forberedelsen. Kant saa foreløbigt i den
nøjagtigere Analyse det vigtigste Middel til Fremgang. Og --- hvad der
for at fastholde Kontinuiteten i hans Udviklingsgang er af særlig
Interesse --- han opgav egentlig aldrig, heller ikke i «Kritik der reinen
Vernunft», det Forhold mellem Analyse og Konstruktion, som han havde
fremstillet i Skriftet «om Grundsætningernes Tydelighed».

Ved sin Fremhæven af Analysens Betydning mødes Kant med Lamberts
Bestræbelser. Men medens Kant ser Konstruktionen ligge i det ubestemte
Fjerne og mener, at det foreløbigt er Analysens Tid, antager Lambert, at
Maalet ligger nærmere. Ved Sammenligning og Kombination af de simple
Begreber, til hvilke Analysen har ført, vil han uden videre danne
Axiomer og Postulater (Brev til Kant, November 1765). Han indser vel, at
en nærmere Prøvelse af de forskellige Kombinationer er
nødvendig\footnote{J.\ H.\ Lambert: \textit{Neues Organon.} Leipzig
1764.\ I, p.\ 323 f.; 499.}; men han føler ikke Trang til at finde særlige
Principer, efter hvilke denne Prøvelse af de mulige Grundsætninger kunde
finde Sted. Han nærmer sig dertil i den mærkelige Udtalelse (Brev til
Kant, Februar 1766), i hvilken han opkaster det Spørgsmaal, om
Erkendelsen af vor Videns Form fører til Erkendelse af vor Videns Indhold.
Men den principielle Vanskelighed ved Overgangen fra Analyse af
Grundbegreberne til Konstruktion af Grundsætningerne er han ikke stanset
ved, og han førte heller ikke Analysen i den bestemte Retning, ved hvilken
Kant lededes til at finde den eneste Mulighed for konstruktiv Erkendelse i
Filosofien. Skønt Kant derfor ærede Lambert som sin Forgænger saa højt,
at han vilde have dediceret ham «Kritik der reinen Vernunft», hvis han
ikke var død før Værkets Fuldendelse\footnote{Se Udkastet til en
Dedikation: \textit{Reflexionen Kants.} II.\ p.\ 1 f.}, saa var det dog
hans Dom om ham, at han bevægede sig inden for den blotte Analyse, uden at
naa det «kritiske» Standpunkt\footnote{\textit{Reflexionen Kants.} II.\
p.\ 67 f.\ (Nr.\ 225, 227, 231).}.

Den Analyse, der spiller saa stor en Rolle i Kants definitive
Erkendelsesteori, angaar den Maade, paa hvilken selve Erkendelsesevnen
virker. Denne Retning fik Kants Analyse allerede, da en omhyggeligere
Bestemmelse af Erkendelsens Grænser viste sig at være nødvendig (se §\
9). Men den førte først til afgørende Resultater, da han fra 1769 af
begyndte at skelne mellem Erkendelsens Stof og Form. Hvorledes dette
Vendepunkt fremkom, skal senere drøftes; her skal det blot undersøges, af
hvad Art den Analyse var, ved hvilken Kant begrundede denne for hans
definitive System saa vigtige Distinktion.

Kant har intet Steds metodisk udført denne Analyse. Han forudsætter den
stedse eller antyder den i stor Korthed. I «Dissertationen» hedder det,
at vi ved at lægge Mærke til den Maade, paa hvilken Bevidstheden virker
paa Erfaringens Foranledning (\textit{attendendo ad ejus actiones
occasione experientiæ}), opdage de Love, efter hvilke denne Virken gaar
for sig. Naar der nu ved den sanselige Erkendelse skelnes mellem Stof og
Form, saaledes at Formen betyder Loven for Bevidsthedens Ordning af det
modtagne Stof, saa fremgaar det af de Udtryk, Kant bruger, at Formen er
det Konstante, Uforanderlige i Sansningen, medens Stoffet er det, der
varierer i det Uendelige. Analysen finder altsaa en Forskel mellem et
Konstant og et Skiftende, et Bestemt og et
Ubestemt\footnote{\textit{Tempus est \ldots\ subjectiva conditio per
naturam mentis humanæ necessaria, \ldots\ quælibet sensibilia certa lege
sibi coordinandi} (Diss.\ §\ 14, 5). --- \textit{Spatium est \ldots\
subjectivum et ideale e natura mentis stabili lege proficiscens, veluti
schema, omnia omnino externe sensa sibi coordinandi.} (Diss.\ §\ 15
D).}. Formen staar som uforanderligt Forbillede (\textit{typus
immutabilis}), der stedse efterlignes paa ny, hver Gang et nyt Indhold
skal ordnes. I Optegnelser fra 70'erne findes det samme Træk.
«Expositionen af det, som tænkes, beror blot paa Bevidstheden; men
Expositionen af det, som er givet, beror --- naar man betragter Stoffet
som ubestemt --- paa Grunden til al Relation og Forbindelse mellem
Forestillingerne». Og naar vi gøre os klart, efter hvilke Love denne
Forbindelse gaar for sig, faa vi «en Original til alle Objekter», «et
Urbillede for enhver mulig Syntese», som kan erkendes afset fra eller
«uafhængigt af den bestemte Sammensætning, i hvilken det var givet»
(\textit{unabhängig von der Einzelnheit darin es gegeben
war})\footnote{\textit{Lose Blätter.} I.\ p.\ 16; 19--21. (Smlgn.\
ovenfor §\ 5 Slutning).}.

Ogsaa i «Kritik der reinen Vernunft» findes Spor af den Analyse, ved
hvilken Formen skilles fra Stoffet. Naar Kant --- ved en uheldig
Udtryksmaade --- siger, at vi kunne erkende Formen forud for virkelig
Iagttagelse, forudsætter han, at en saadan Analyse allerede er foretagen,
og ved virkelig Iagttagelse mener han: enhver bestemt, særskilt
Iagttagelse. --- Nogle Exempler skulle vise dette.

«Naar jeg fra Forestillingen om et Legeme udsondrer hvad Forstanden
tænker derom (Substants, Kraft, Delelighed o.\ s.\ v.), og ligeledes det
deraf, som hører til Sansefornemmelsen (Uigennemtrængelighed, Haardhed,
Farve o.\ s.\ v.), saa faar jeg endnu Noget tilbage af denne empiriske
Anskuelse, nemlig Udstrækning og Form. Disse høre til den rene Anskuelse,
der finder Sted i Bevidstheden ogsaa uden nogen virkelig Genstand for
Sansning eller Fornemmelse, som Sansningens blotte Form». Der beskrives
her en Analyse, d.\ v.\ s.\ en successiv Henvendelse af Opmærksomheden
paa det Givnes forskellige Sider eller Elementer, og Analysens Resultater
deles i tre forskellige Klasser, der henføres til de tre Evner, Forstand,
Anskuelsesform og Sansefornemmelse\footnote{Kant udtaler sig her paa den
gamle Abstraktionsteoris Maade, skønt han andre Steder bestemt forkaster
denne. Se «Kritik der reinen Vernunft». 2.\ Aufl.\ p.\ 457 (Noten til
første Antithesis). \textit{Reflexionen Kants.} II.\ p.\ 126 (Nr.\
412).}. Hvorledes «Udsondringen» nærmere begrundes, ses af andre Steder.
«Rum og Tid ere [Sansningens] rene Former, Fornemmelse over Hovedet er
Stoffet. \ldots\ Hine hænge aldeles nødvendigt sammen med vor Sansning,
af hvilken Art saa vore Fornemmelser ere; disse kunne være meget
forskellige». «Vi kalde Syntesen af det Mangfoldige i Indbildningskraften
transscendental, naar den afset fra enhver Forskel i det Anskuede
(\textit{ohne Unterschied der Anschauungen}) ikke gaar paa Andet end
Forbindelsen af det Mangfoldige a priori». «Apperceptionen og med den
Tænkningen gaar forud for al mulig bestemt Anordning af
Forestillingerne»\footnote{\textit{Kritik der reinen Vernunft.} 2.\
Aufl.\ p.\ 35; 60; 118; 345.}. --- Og paa et Sted i «Prolegomena», hvor
Kant beretter om Kategorilærens Tilblivelse, siger han: «Jeg saa mig om
efter en Forstandshandling, der indeholdt enhver anden og kun frembød
Forskelligheder ved de forskellige Modifikationer eller Momenter, under
hvilke Forestillingens Mangfoldighed bringes ind under Tænkningens Enhed.
Da fandt jeg, at denne Forstandshandling bestod i at
dømme»\footnote{\textit{Prolegomena.} Riga 1783.\ p.\ 109.}.

Der stillede sig tillige den Opgave for Kant at finde fælles Kendemærker
for de konstante Elementer, han saaledes efter Haanden kunde udskille i
vor Erkendelse. Denne Del af Analysen gik Haand i Haand med Paavisningen
af de konstante Elementer, idet Kant først blev ret opmærksom paa disse,
da han havde fundet en fælles Grundbestemmelse for al vor Erkendelse.
Denne Grundbestemmelse var Forbindelsen, Syntesen, og da den er en
Funktion, ansaa Kant det som selvfølgeligt, at den ikke kunde være given
med de blotte Indtryk. Men om denne Side af Analysen skal der handles
nærmere i Kap.\ 4. Her holde vi os blot til den almindelige Betydning,
den analytiske Metode vedblev at have for Kant.

Denne Analyse, der gaar ud paa at finde selve Erkendelsesevnens Maade at
virke paa, skelnes udtrykkeligt fra «den sædvanlige Fremgangsmaade i
filosofiske Undersøgelser: at analysere de sig frembydende Begreber efter
deres Indhold og bringe dem til Tydelighed». Den kritiske Filosofis
egentlige Opgave er en Analyse af selve Erkendelsesevnen, at finde
Grundbegreberne dér, hvor de udspringe, og paavise deres Anvendelse. Og
denne kritiske Analyse, der skiller Form fra Stof, kan ikke gaa for sig,
«før lang Øvelse har gjort os opmærksom og sat os i Stand til at udsondre
det, som vor egen Erkendelsesevne yder, fra det, vi modtage ved
Sanseindtryk». Derfor har Kant heller ikke Noget imod at indrømme
Matematikeren Kästner, at Rummets Begreb er abstraheret af sanselige
Forestillinger: «ti uden Anvendelse af vor sanselige Forestillingsmaade
paa virkelige Genstande for Sanserne, vilde selv det, der a priori
indeholdes i disse [o:\ i Sanserne], slet ikke blive os
bekendt»\footnote{Dissert.\ §\ 15 Coroll. --- \textit{Kritik der reinen
Vernunft.} 2.\ Aufl.\ p.\ 1 f.; 90. --- Wilh.\ Dilthey: \textit{Aus den
Rostocker Kanthandschriften. Archiv für die Geschichte der Philosophie.}
III.\ p.\ 87.}.

13.\ Det ved Analyse fundne Grundlag for den kritiske Erkendelseslære
træder stærkere frem i førsteudgave af «Kritik der reinen Vernunft» end i
Prolegomena og anden Udgave\footnote{Dette er paavist af Benno Erdmann:
\textit{Einleitung zu Prolegomena,} p.\ XXXII--XXXVIII.}. Hist udvikler
han paa en i psykologisk Henseende meget interessant Maade Begrebet om
Syntesen som den almindelige Form for Bevidsthedsvirksomheden paa de
forskellige Trin. Senere lader han sig nøje med ganske i Almindelighed at
henvise til den forbindende Proces, særligt paa Forstandens
Omraade\footnote{Kfr.\ ogsaa Brevet til Tieftrunk af 11.\ December 1797 om
Begrebet Sammensætten som det, der ligger til Grund for alle
Kategorier.}. Paa Grund af de Misforstaaelser i psykologisk-idealistisk
Retning, «Kritik der reinen Vernunft» strax blev Genstand for, søgte Kant
i de senere Bearbejdelser at bortlede Opmærksomheden fra den psykologiske
Analyse af Erkendelsesprocessen (fra det, som han i Fortalen til første
Udgave kaldte den subjektive Deduktion), der kun havde en forberedende
Betydning, for at hans Hovedopgave, Konstruktionen af de Fornuftsætninger,
der gælde for alle Fænomener og dog ere uafhængige af disse, kunde træde
des skarpere frem.

Men den skarpe Modsætning mellem Matematikens og Filosofiens Metode, som
Kant havde fremstillet i Afhandlingen om Grundsætningernes Tydelighed,
fastholder han dog stedse. I det interessante Afsnit «Die Disciplin der
reinen Vernunft im dogmatischen Gebrauche» udtales den til Dels med de
samme Vendinger som i det næsten tyve Aar ældre Skrift. Ogsaa her vises
det, at de filosofiske Begreber, der vindes ved Analyse, ikke som de
matematiske ved direkte Konstruktion, ikke i strengeste Forstand kunne
defineres, da der ikke kan haves nogen apodiktisk Sikkerhed for, at
Begrebets Indhold er fuldstændigt. Som Exempler nævnes Begreberne Aarsag,
Ret og Billighed. Der kan være dunkle Biforestillinger til Stede, som vi
ikke blive opmærksomme paa ved Analysen, og som dog uden at vi mærke det
spille en Rolle med ved Anvendelsen\footnote{\textit{Kritik der reinen
Vernunft.} 2.\ Aufl.\ p.\ 756. --- Smlgn.\ hermed «Träume eines
Geistersehers». 1.\ Th.\ 1.\ Hauptst.\ (Kehrbachs Udgave p.\ 6 Note),
hvor Beskrivelsen af tilsnegne Begreber ganske minder om det anførte Sted
af «Kritiken».}.

Forskellen mellem de ældre Skrifter (1762 ff.) og «Kritik der reinen
Vernunft» er, at medens Kant tidligere skelnede mellem to Slags Begreber,
empiriske og matematiske, af hvilke hine vindes ved Analyse, disse ved
direkte Konstruktion, skelner han senere tre Slags Begreber: 1)
empiriske, 2) apriorisk-diskursive, 3) apriorisk-intuitive. De to første
Arter vindes ved Analyse, den sidste ved direkte Konstruktion. Men den
Analyse, ved hvilken den første og den anden Art vindes, er, som det er
paavist i det Foregaaende (se ovenfor §\ 12), af forskellig Beskaffenhed.

14.\ For Kritiken af Kants Erkendelsesteori er denne Sammenhæng mellem
det ældre og det senere Standpunkt af stor Interesse. Ti den lægger det
Spørgsmaal nær, om Kant virkeligt har opdaget en helt ny Klasse af
Begreber, der skulde staa imellem de empiriske og de konstruerede. Ved
Afgørelsen af dette Spørgsmaal behøver man ikke at gaa ud over Kants egne
Udsagn. Ti han erklærer jo med rene Ord, at de apriorisk-diskursive
Begreber (o:\ Kategorierne) ikke ere fuldt færdige --- og dog opererer
han med dem til Opstilling af Grundsætninger, som om de vare fuldt
færdige! Se vi hen til deres Tilblivelsesmaade, saa ere de samme Vilkaar
underlagte som de sædvanlige empiriske Begreber. De skyldes, som vi have
set, Iagttagelsen af konstante Elementer under vore Erfaringers Skiften.
Antagelsen af, at vi her have selve Erkendelsesevnens Natur for os
(«Kilderne», som Kant plejer at sige), er en Slutning fra eller en
Fortolkning af denne Konstants. Den er altsaa en Hypotese. I Fortalen til
første Udgave af Fornuftkritiken siger Kant ogsaa, at «den subjektive
Deduktion» (o:\ Paavisningen af Erkendelsens Kilder formedelst Analyse)
kunde se ud som en Hypotese, da den jo gaar ud paa «Paavisningen af en
Aarsag til en given Virkning»; men han mener, at det ikke er saa, og
lover ved en anden Lejlighed at godtgøre det. Dette Løfte har han
imidlertid ikke holdt, hvorimod han i sine senere Fremstillinger har
reduceret dette analytiske Parti betydeligt.

Kant er altsaa i en vis Forstand paa sit definitive Standpunkt ikke naat
videre end han var 1762. Han har ikke paavist Muligheden af at bestemme
«Formerne» saaledes, at der er fuld Garanti for Bestemmelsens Nøjagtighed
og Fuldstændighed. Han kan da egentlig ikke vide sig sikker mod den samme
Anke, som han rettede mod Leibniz og Wolf: at operere med ufærdige og
tilsnegne Begreber. Og der er et betydeligt hypotetisk Element i det
System af Former, han stiller op som Fundament for vor exakte Erfaring.
Han har Ret i, at hvis vi kunde være ganske sikre paa, at disse Former
maa spille en Rolle med ved al vor Erfaring, al vor Opfattelse, saa vilde
vi ogsaa have en apriorisk Viden om Erfaringsgenstandene. Men netop denne
fulde Sikkerhed kan ikke godtgøres.

For Erkendelsesteorien er der her ingen anden Vej at gaa end at operere
med det ved Analyse fundne Grundlag som med en Hypotese, idet man
overlader til fortsat Analyse at berigtige det Grundlag, hvorfra man gaar
ud. Sagen er, at de forskellige Operationer staa i Vexelforhold til
hinanden. Man kan ikke lægge al mulig Konstruktion bort og give sig til
at analysere; det vil være til Skade for selve Analysen, da
Konstruktionen kan være nødvendig som Hjælpeproces og Experiment. Kants
egen antinomiske Metode var egentlig Forsøg paa Konstruktioner, der
skulde belyse Forudsætningernes Beskaffenhed, naar disse ikke direkte
kunde analyseres. --- Ogsaa det Omvendte gælder naturligvis, at
Konstruktionen forudsætter Analysen, hvorfor selv de matematiske
Begreber, der synes opstillede ved direkte Konstruktion, ved nærmere
Undersøgelse vise sig at forudsætte en foregaaende Analyse. Kants
Modsætning mellem Matematik og Filosofi er altsaa ikke holdbar i den
Form, i hvilken han har opstillet dem.

Visheden om, at vi i Erfaringens konstante Elementer have vor
Bevidsthedsvirksomheds nødvendige Former, kan, saa langt end Analysen
føres, kun blive approximativ, og de Slutninger, der kunne bygges paa
denne Vished, maa naturligvis dele Skæbne med
den\footnote{Smlgn.\ allerede \textit{Ænesidemus} [af G.\ E.\ Schulze].
1792 p.\ 401; 406.}. Kants Tendents paa hans definitive Standpunkt til at
skyde Analysen til Side har gjort ham til Dogmatiker mod hans Villie. Han
saa rigtigere i 1762, da han brugte den analytiske Metode som Vaaben mod
Dogmatikerne.

15.\ At Analysen har været Stedbarn i Kants definitive Filosofi, skønt
den dog har sin bestemte Plads i denne, viser sig især i, at Udledningen
af «Formerne» af den konkrete Sammenhæng, i hvilken de efter Kant fra
først af lægge sig for Dagen, ikke er gennemført. Forholdet mellem de to
Sætninger, som Kant begge vedkender sig: «Al Erkendelse begynder med
Erfaring», og: «Ikke al Erkendelse udspringer af Erfaring», har han ikke
tydeligt paavist. Dette vilde kun være sket, hvis han havde fremstillet
de forskellige Udviklingstrin, «Formerne» gennemløbe fra det medfødte
Grundlag til den ved Øvelse erhvervede bevidste
Klarhed\footnote{Smlgn.\ Distinktionen mellem Medfødthed og oprindelig
Erhvervelse i Dissert.\ §\ 8; Coroll. --- \textit{Ueber eine Entdeckung}
etc.\ p.\ 68. --- \textit{Kritik der reinen Vernunft.} 2.\ Aufl.\ §\ 27.
--- Brev til Hertz af 21.\ Febr.\ 1772. --- Fra et rent psykologisk
Standpunkt maa der ogsaa opponeres mod Kants Distinktion mellem Stof og
Form. Fornemmelserne ere ikke blot passivt optaget Stof, men hver
Fornemmelse bestemmes i sin Opstaaen og Beskaffenhed ved Bevidsthedens
hele Tilstand. Her finder for saa vidt, hvad selve Fornemmelserne angaa,
en Formen og Forbinden, en Syntese Sted. Smlgn.\ min Psykologi 3.\ Udg.\
p.\ 130. Denne Indvending ophæver ikke Kants Distinktion, men gennemfører
den videre og paa anden Maade end han gjorde.}.

Men ikke blot forudsætter Kant Formerne som fuldtfærdige; han forudsætter
dem ogsaa i en ideal Fuldkommenhed og Renhed, i hvilken de ikke forekomme
i nogen virkelig Anskuelse eller Erfaring. At dette er saa for Rummets og
Tidens Vedkommende, behøver ingen nærmere Paavisning. Newtons absolute
Rum og Tid ere hos Kant saa at sige slaaede ind, ere blevne til subjektive
Former --- uden dog at miste deres Absoluthed. Saaledes betegnes enhver
af Anskuelsesformerne i Dissertationen som \textit{typus immutabilis} (§\
16 Coroll.). Kant tænker sig altsaa Forbilleder, ideale Rammer, til
hvilke Alt, hvad der fremtræder for os, maa henføres for at indvirkes i
en almindelig Sammenhæng\footnote{Hvad Tiden angaar, har Kant ikke
skelnet mellem Tidsfornemmelse, Tidsforestilling og Tidsanskuelse.
Smlgn.\ min Psykologi.\ 3.\ Udg.\ p.\ 206--214 og min Afhandling: «Lotze
og den svenske Filosofi». \textit{Nordisk Tidsskrift,} udg.\ af den
Letterstedtske Forening 1890.\ p.\ 6--10. --- Hvad Rummet angaar,
udtrykker Kant sig i andre Skrifter langt tydeligere end i «Kritik der
reinen Vernunft». Det hedder saaledes i «Metaphysische Anfangsgründe der
Naturwissenschaft» (Riga 1786).\ p.\ 3 f.: «Der absolute Raum ist an sich
nichts und gar kein Object, sondern bedeutet nur einen jeden andern
relativen Raum, den ich mir äusser dem gegebenen jederzeit denken kann».
--- p.\ 16: «Der absolute Raum ist für alle mögliche Erfahrung
nichts».}. Et lignende Forbillede (\textit{exemplar}) dannes i Følge
Dissertationen (§\ 9) ogsaa for Forstandserkendelsens Vedkommende og
afgiver den fælles Maalestok for al Realitet. I «Kritik der reinen
Vernunft» opstaar selve Erfaringsbegrebet eller Naturbegrebet ved en
lignende ideal Konstruktion. Parallelen med hint «uforanderlige
Forbillede» træder tydeligt frem, naar det hedder: «Der gives kun én
Erfaring, i hvilken alle Iagttagelser forestilles i gennemgaaende og
lovmæssig Sammenhæng, ligesom der kun er ét Rum og én Tid; i hin finder
alle Fænomener og alle Forhold af Væren og Ikke-Væren
Sted»\footnote{\textit{Kritik der reinen Vernunft.} 1.\ Aufl.\ p.\ 110.}.
«Alle vore Erkendelser ligge inden for al mulig Totalitet». «Alle
Fænomener ligge i én Natur og maa ligge deri, fordi der uden den
aprioriske Enhed ingen Enhed af Erfaringen og altsaa ingen Bestemmelse af
Genstandene i den er mulig»\footnote{\textit{Kritik der reinen Vernunft.}
2.\ Aufl.\ p.\ 185; 263.}. --- Den sidste Begrundelse, Kant giver af
denne Erfaringens (eller Naturens) Enhed, er, at bevidst Erkendelse kun
er mulig, naar de givne Elementer forenes ved Syntese. Den syntetiske
Enhed er Princip for al Forstaaelse. «Naar enhver enkelt Forestilling var
ganske fremmed for, isoleret og sondret fra de andre, saa vilde aldrig
noget Saadant, som Erkendelse er, kunne opstaa. Ti Erkendelse er et Hele
af sammenlignede og forbundne Forestillinger»\footnote{\textit{Kritik der
reinen Vernunft.} 1.\ Aufl.\ p.\ 97 (kfr.\ 2.\ Aufl.\ p.\ 103).}. Vi
have derfor, som vi allerede have vist (§\ 5 Slutn.; §\ 12), i Begrebet
om vor erkendende Bevidstheds syntetiske Karakter et Urbillede for Alt,
hvad vi kunne erkende.

Men i hvilken Grad tilfredsstille nu de virkelige Iagttagelser den ideale
Fordring, der ud fra dette Forbillede stilles til dem? Derom kan
aabenbart Intet vides paa Forhaand. Et ved Abstraktion og Idealisation af
vor Erkendelsesvirksomhed dannet Forbillede viser sin Betydning ved at
sætte Tanken i Bevægelse. Men om Tanken finder adækvate Genbilleder af
Forbilledet, og om den kan gennemføre Forbilledet midt i Iagttagelsens
uoverskuelige Mangfoldighed, det er det store Spørgsmaal, som ikke er
besvaret med Paavisningen af, at vi kun have nødvendig Erkendelse, hvis
det sker.

Undertiden udtaler Kant sig da ogsaa, som om der ved Forbilledet blot
indledtes en Søgen. I Dissertationen opføres saaledes det Princip, at Alt
hvad der sker, sker efter Naturens Orden, som et subjektivt Princip for
vor Forsken (\textit{principium convenientiæ}), paa Linie med
Sparsommelighedens Lov (§\ 30). I en Optegnelse fra 70'erne tales om en
Præsumtion, i Følge hvilken Alt kan bestemmes efter en Regel. I en
Optegnelse «fra Kriticismens første Tid» spørges: «Da vi ikke indse
Muligheden af en Realgrund, hvorledes kunne vi da a priori indse, at der
absolut maa gives saadanne? \ldots\ Gælder dette Princip ikke som
Anticipation, fordi vi uden Regel ingen Erfaringer vilde have, og denne
Regel bestaar i Tidens og Rummets Ordning efter almindelige Love?» I en
Optegnelse fra 80'erne hedder det: «Principerne for Fænomenernes
Exposition er Principer for deres Intellektion, ikke for deres
Perspicients, for at søge Aarsager, ikke for at bestemme dem». Og i
«Kritik der reinen Vernunft»: «Forstanden er stedse beskæftiget med at
gennemspejde Fænomenerne i den Hensigt at udfinde en Regel for dem».
«Forstanden kan a priori ikke yde Mere end at anticipere en mulig
Erfarings Form over Hovedet».\footnote{\textit{Lose Blätter.} I.\ p.\ 35;
136. --- \textit{Reflexionen Kants.} II.\ p.\ 307 (Nr.\ 1071). ---
\textit{Kritik der reinen Vernunft} 1.\ Aufl.\ p.\ 126; 2.\ Aufl.\ p.\
303.}

Hvis Kant var bleven staaende herved, vilde han ikke have paastaaet Mere
end hvad han kunde bevise. Men da vilde han ganske vist ikke have havt
nogen Grund til at opstille sin skarpe Forskel mellem Kategori og Ide,
eftersom det da vilde vise sig, at allerede Kategorien stiller ideale
Fordringer, som den reale Erkendelse stedse kun ufuldkomment
fyldestgør\footnote{Allerede Salomon Maimón paaviste dette. «Für mich ist
Erfahrung im strengen Sinne kein in der Anschauung darstellbarer Begriff,
sondern eine Idee». \textit{Kritische Untersuchungen.} Leipzig 1797.\ p.\
154. --- «Die kritische Philosophie kann nur hypothetisch philosophiren».
\textit{Die Kategorien des Aristoteles.} Berlin 1794.\ p.\ 133.}. Kant
blev ført ud over Maalet ved sin Tillid til den «kopernikanske» Tanke,
hvori (som senere skal omtales) Opdagelsen 1769 bestod. Den gav ham Mod
til at foretage den syntetiske Konstruktion, der endnu faa Aar forud stod
for ham som et fjernt Maal. Fra Urbilledet og Præsumtionen gik han uden
videre til Loven, som om denne kunde udledes af hine. Han anser den
virkelige Erkendelse for given med Erkendelsens Begreb: den gamle
ontologiske Tankegang har her overlistet sin store Modstander! Den
transscendentale Deduktion, som havde kostet ham saa meget Besvær, og som
han var saa stolt over, var i Virkeligheden erkendelsesteoretisk Alkymi.
Men som Alkymisterne ofte i deres Søgen efter de Vises Sten under Vejs
opdagede nye kemiske Forhold, saaledes kom Kant under sin ivrige Stræben
efter at give en absolut Løsning af Humes Problem til at give værdifulde
Bidrag til Forstaaelse af vor Erkendelses Natur.

16.\ Kun et enkelt Punkt skal endnu fremdrages til Belysning af Kants
analytiske Metode. Han gaar, som vi have set (§\ 12), ud fra, at de
konstante Elementer i vor Erfaring maa skrive sig blot fra selve vor
Erkendelsesvirksomhed og kun være Udtryk for dennes Love. Dersom Kant
udførligere var gaaet ind paa den Analyse, der førte ham til Opstilling
af Formerne, maatte han ogsaa være kommen til nøjere at prøve
Berettigelsen af denne Antagelse. Han gaar jo fra første Færd ud fra, at
Bevidstheden, det erkendende Subjekt ikke er det eneste Existerende, men
kun er et Led af den hele Tilværelse og staar i Vexelforhold med det
øvrige Existerende. Den Maade, hvorpaa Subjektet opfatter Tilværelsen,
kan da ikke blot være bestemt ved dets egen Natur, men maa ogsaa være
bestemt ved det Existerendes Natur.

Der foreligger adskillige, ikke paaagtede Antydninger hos Kant i denne
Retning. I «Kritik der reinen Vernunft», første Udgave (p.\ 358), tales
om «det Noget, der ligger til Grund for de ydre Fænomener og paavirker
vor Sans saaledes, at den faar Forestillingerne om Rum, Materie, Figur
o.\ s.\ v.». Og i Anledning af det Spørgsmaal, hvorledes Rumsanskuelse i
det Hele er mulig i et tænkende Subjekt, bemærkes, at dertil kun kan
siges, at Forestillingerne om ydre Fænomener have deres Aarsag i Tingen i
sig selv («den transscendentale Genstand») (ibid.\ p.\ 393). Den Maade,
jeg opfatter en Ting paa, behøver jo slet ikke at ligne selve den Ting,
der fremkalder Opfattelsen (Sansefornemmelsen eller Rumsanskuelsen) hos
mig --- hedder det i Prolegomena p.\ 64. Rumsanskuelsen er, hedder det i
anden Udgave af Fornuftkritiken (p.\ 66, 69), bestemt ved Genstandens
Forhold til Subjektet. De tydeligste Udtalelser findes paa nogle Steder i
«Dialektiken». Det er umuligt, hedder det, at besvare det Spørgsmaal,
«hvorfor den transscendentale Genstand giver vor Sanseanskuelse netop kun
Rumsanskuelse og ikke en anden Art Anskuelse». «Mange Kræfter i Naturen,
der ytre deres Existents ved visse Virkninger, ere og blive uudforskelige
for os; ti vi kunne ikke efterspore dem vidt nok ved Iagttagelse. Det til
Grund for Fænomenerne liggende transscendentale Objekt, og med det
Grunden til, at vor Sansning netop har disse bestemte øverste Betingelser,
og ikke andre, ere og blive uudforskelige for os, skønt Sagen selv
iøvrigt er given, blot ikke gennemskuet» (2.\ Udg.\ p.\ 585; 631 f.). ---
Paa disse sidste Steder bliver Grunden til Rumsanskuelsen endog aldeles
overvejende tilskreven Tingen i sig.

Af disse Antydninger ses, at Kant ikke blot udledte Stoffet, men ogsaa
Formen fra Tingen i sig. Men naar saa er, er det uberettiget, naar han
saa ofte opstiller et Dilemma, der giver Valget mellem at lade Genstanden
bestemme Forestillingen og Forestillingen bestemme Genstanden. I hans
egen Teori er i Virkeligheden begge Dele forudsat. Kant kan derfor ikke
konsekvent være den strenge Aprioriker, han vil være. Der er en
Grundforudsætning, han ikke har draget frem, og paa hvilken dog hele hans
«Grundsætningernes System» hviler: nemlig den, at Tingen i sig virker
konstant. Er det nemlig ikke Tilfældet, saa vilde jo, selv om Subjektets
Natur forudsættes uforanderlig, de konstante Elementer i Erfaringen, der
ere det virkelige Grundlag for Kants hele erkendelsesteoretiske Bygning,
kunne ændres, og Grundsætningerne vilde da ogsaa ændres. --- Ogsaa fra
denne Side viser det sig, at «den kritiske Filosofi kun kan filosofere
hypotetisk». Kants Antagelse, at han kunde gennemføre en Filosofi om
Fænomenerne, uden at behøve at opstille nogen som helst Forudsætning om
Tingen i sig selv, var urigtig. Hans første Kritikere gjorde allerede
opmærksom paa, at naar han bruger Begreberne Ting, Existents og Aarsag om
den, er den dog ikke absolut ubekendt. Men det maa føjes til, at han
ogsaa forudsætter, at den virker paa en regelmæssig Maade.

Kant mente at have ført et Bevis for Aarsagssætningens Gyldighed, Noget,
som Wolf var kommen højst uheldigt fra. Medens Wolf vilde udlede den af
Modsigelsens Grundsætning, begrunder Kant den som Forudsætning for
Erfaringens Mulighed\footnote{\textit{Kritik der reinen Vernunft.} 2.\
Aufl.\ p.\ 264.}. Kants Bevis hviler paa den Antagelse, at
Aarsagsbegrebet er en af de subjektive Former, uden hvilke Erfaring i
streng Betydning ikke bliver til. Men lige saa lidt som den
Omstændighed, at vi opfatte ydre Fænomener i Rummets Form, kunde udledes
af Subjektets Natur alene, lige saa lidt anser Kant det for at være
uafhængigt af Tingen i sig, at vi opfatte Fænomenerne som Aarsager og
Virkninger. Dette ses af en Udtalelse som denne: «Dem transscendentalen
Objecte können wir allen Umfang und Zusammenhang unserer möglichen
Wahrnehmungen zuschreiben \ldots\ Die Erscheinungen sind, ihm gemäss,
\ldots\ gegeben».\footnote{Ibid.\ p.\ 522 f.} Her siges altsaa
Fænomenerne ikke blot hvad Stof, men ogsaa hvad Formen angaar, at være i
Overensstemmelse med Tingen i sig, idet ogsaa deres Sammenhæng henføres
til den. Tingen i sig maa altsaa være lige saa konstant i sin Virken paa
Subjektet, som Sammenhængen mellem Fænomenerne er. Og her staa vi da ved
en Grundforudsætning om Tingen i sig, som Kant aldrig har formuleret. Det
viser sig, at der til «Erfaringens Mulighed», der bærer hele Kants
Erkendelsesteori, hører Mere end de subjektive Betingelser. Herved
indskrænkes baade Apriorismen og Fænomenalismen. Muligheden for at
antecipere Fænomenerne og Nødvendigheden af at skelne mellem disse og
Tingen i sig begrænses paa betydningsfuld Maade. Navnligt bortfalder
Nødvendigheden af at udelukke Aarsagsprincipet fra Gyldighed for Tingen i
sig selv. Skønt det strenge, aprioriske Bevis for Aarsagssætningen
bortfalder, opnaas der den Fordel, at den store Vanskelighed, som laa i
Kants Antagelse af Tingen i sig som Aarsag til Fornemmelserne, falder
bort med det Samme. Samtidigt med, at Systemet ændres i empirisk og
realistisk Retning, mister det ogsaa sin store Inkonsekvents.

\section*{III.\\ Teori og Praxis.}

\noindent 17.\ Vi vende paa ny tilbage til Kants Standpunkt i 1762 og de
følgende Aar, for at betragte det fra en ny Side. Kant havde indset
Nødvendigheden af en Reformation af Filosofien og ventede foreløbigt
Frelsen ad Analysens Vej. Men naar saa Meget stod uafgjort og ufærdigt
hen, og naar vi med Hensyn til nogle af Grundbegreberne (f.\ Ex.\
Aarsagsbegrebet) stode over for det Uforklarlige, saa kan det Indbegreb
af teoretisk Filosofi, Kant paa denne Tid havde til Raadighed for positiv
Fremstilling, ikke have været stort. I sin Bestræbelse for «hellere at
have en grundig Filosofi end en vidtløftig» havde han gjort rent Bord.
Selv i «Beweisgrund», det mest konstruktive Skrift fra denne Tid,
erklærer han, at Meningen ikke var at give en egentlig Bevisførelse, og
at Guds Tilværelse heller ikke behøvede at bevises. I Tillid til, at den
religiøse Tro havde et helt andet Grundlag end Bevisførelse, omstyrtede
han den naturlige Teologis hele Basis. Og paa lignende Maade gik han i
«Träume» frem over for den spiritualistiske Psykologi. De kosmologiske
Antinomier synes samtidigt at have spillet en stor Rolle for ham ved hans
Forsøg paa at finde Erkendelsens Grænser. Materialet til det Parti af
«Kritik der reinen Vernunft», som han kaldte Dialektik, og som var det,
der ikke gjorde mindst Virkning ved Værkets Fremkomst, var saaledes i det
Væsentlige færdigt.

Kants Filosofi maatte derved faa et vist negativt Præg i Modsætning til
den af Positivitet struttende Dogmatisme. Kant var sig dette fuldt
bevidst. Han bruger om sin Filosofi Udtrykket «Uvidenhedens Filosofi»,
ogsaa «negativ Filosofi». Han finder denne Filosofi den vanskeligste af
alle, fordi den maa gaa tilbage til Erkendelsens Kilder for at kunne
begrunde sin Ikke-Viden.

«Hvor Vildfarelse er fristende og tillige farlig, dér ere negative
Erkendelser og Kriterier vigtigere end positive og udgøre ofte vor
Videnskabs egentlige Objekt \ldots\ Sokrates havde en negativ Filosofi,
hvad Spekulationen angik, nemlig om megen formentlig Videnskabs
Værdiløshed og om Grænserne for vor Viden. Den negative Del af
Opdragelsen er den vigtigste
(Rousseau).»\footnote{\textit{Reflexionen Kants.} II.\ p.\ 44 f.\ (Nr.\
148--151).}

Disse Henvisninger til Sokrates og til Rousseau ere karakteristiske for
Kants Standpunkt paa denne Tid. Sokrates's Viden var jo en Viden om, at
han Intet vidste, en Viden, der var begrundet paa klar Indsigt i hvad der
hørte til for at vide Noget. Og Sokrates advarede jo ogsaa mod ufærdige
Begreber og søgte ved alsidig Drøftelse at føre til rigtige
Begrebsbestemmelser. Hvad Rousseau angaar, drager Kant i den anførte
Udtalelse en Parallel mellem hans «negative Opdragelse» og sin egen
negative Filosofi. Ved negativ Opdragelse forstod Rousseau «en
Opdragelse, der stræber at fuldkommengøre Organerne, vor Erkendelses
Instrumenter, før den giver os Erkendelsen selv, og som forbereder til
Fornuft ved at øve Sanserne. Den negative Opdragelse lægger ikke Hænderne
i Skødet, langt fra: den bibringer ingen Dyder, men den forebygger
Laster; den lærer ikke Sandheden, men den forebygger Vildfarelser.» Den
begrundes særligt ved Nødvendigheden af at lære Individets Natur og
Karakter at kende, saaledes som de uvilkaarligt udfolde sig, før man
griber ind i Henhold til almindelige og overleverede Regler. Den er i sin
Omhu for at holde Vildfarelse ude og for at finde den ejendommelige
Naturs Grænser en vanskelig Kunst --- den Kunst: \textit{tout faire en ne
faisant rien.}\footnote{\textit{Lettre à M.\ de Beaumont} (Petits chefs
d'œuvres de J.\ J.\ Rousseau.\ Paris 1859.\ p.\ 313). --- \textit{Émile.}
Livre II.\ (Paris 1851.\ p.\ 80 f.; 116.)} --- I flere Henseender har den
kritiske Filosofi efter det Bedste og mest Berettigede i dens Aand sit
Forbillede i Rousseau's Lære om den negative Opdragelse. Begge Steder
tjene Negation og Kritik til at værne om Livets positive Udfoldelse,
medens Dogmatismen, baade i Filosofien og i Pædagogiken, er ivrig for at
gribe ind og fastsætte eller regulere Alt efter sine «evige Sandheder».

Det er bekendt, i hvor høj Grad Rousseau's Skrifter interesserede Kant.
Da han fik Emile, holdt den ham tilbage fra hans sædvanlige Spadseretur.
Men hvis Kants Optegnelser og Tilføjelser til hans «Beobachtungen über
das Schöne und Erhabene» ikke i sidste Øjeblik vare blevne reddede ud af
en Urtekræmmers Makulaturbunke\footnote{\textit{Sämmtl.\ Werke.} Ed.\
Rosenkranz u.\ Schubert XI, 1.\ p.\ 218 f.}, vilde man ikke have vidst,
hvor dybt personligt dette Udtryk var. Det hidførte for Kant et helt nyt
Grundlag for Vurdering af Mennesker og menneskelige Forhold. Hidtil havde
Kant været Optimist, havde betragtet den intellektuelle Udvikling som det
Højeste og set Fremskridtet afgjort ved den. Af Rousseau lærte han en
Maalestok for Menneskeværd, der til en vis Grad var uafhængig af den
intellektuelle Udvikling. Han lærte nu, at Pøbelen ikke er at foragte
blot fordi «den Intet véd». Han «lærte at ære Menneskene». Og han priste
Rousseau, fordi han havde fremdraget Menneskets dybt skjulte Natur under
Civilisationens Former, der altfor ofte tildække
dem.\footnote{\textit{Sämmtl.\ Werke.} Ed.\ Rosenkranz u.\ Schubert XI,
1.\ p.\ 240.\ 248.} --- Der var da her givet et faktisk Grundlag for
Menneskelivet, hvorved dette blev uafhængigt af hvorledes det gik med den
videnskabelige Erkendelse af Tilværelsen og dens Ophav. Det var for Kant
en stor Støtte, da han jo netop havde grundlagt sin Kritik af den
naturlige Teologi og af den spiritualistiske Psykologi. Om de Lærde
mistede nogle Yndlingsdoktriner, saa mistede «den store, for os mest
agtelsesværdige Mængde»\footnote{\textit{Kritik der reinen Vernunft.} 2.\
Aufl.\ Vorrede p.\ XXXIII.} Intet.

18.\ Men det forholder sig med denne Rousseau'ske Paavirkning, som med
den Hume'ske: Dersom den ikke var sikret ved ydre Vidnesbyrd, vilde vi i
Kants Skrifter ikke finde nogen tvingende Anledning til at forudsætte
den. I og for sig vilde man meget godt af Kants Tankegang, saaledes som
den udfoldede sig i Aaret 1762 og de følgende Aar, have kunnet forstaa,
hvorledes han kom til den Distinktion mellem Teori og Praxis, der fra
denne Tid af --- altsaa længe før han slog den fast i sine Fornuftkritiker
--- spillede saa stor en Rolle hos ham.

Med Kritiken af Beviserne for Guds Tilværelse maatte ogsaa Dobbeltheden
af de Motiver, der havde ført til Konstruktionen af Gudsbegrebet, komme
for Dagen. Den ene Gruppe af Motiver var af rent teoretisk Art og bestod
væsentligt i Trangen til at afslutte Aarsagsrækken, at lade de arbejdende
Tanker komme til Ro i Kontemplationen af et «i sig selv grundet» Væsen.
Den anden Gruppe var af æstetisk-etisk Art, og bestod væsentligt i en
Trang til i ét Begreb at koncentrere alt Værdifuldt i Verden. Ved
Sammensmeltningen af begge Motivgrupper opstod den dogmatiske Filosofis
Begreb om Gud som \textit{ens realissimum}, Indbegrebet af al Realitet,
saaledes at Realitet betød det Samme som Fuldkommenhed, idet kun saadanne
Egenskaber, som udtrykke Negation og Begrænsning, skulde være udelukkede.
Allerede i «Beweisgrund» (I, 4, 3) opponerede Kant mod denne
Sammenblanding af Begreberne Realitet og Fuldkommenhed. Fuldkommenhed
udtrykker en Vurdering, som foretages af et følende og villende Væsen:
den kan altsaa ikke være en absolut Egenskab. Og i Afhandlingen om de
negative Størrelser (III, 2.\ Anm.) gør Kant opmærksom paa, at da Ulyst
er lige saa positiv og real som Lyst, kan man ikke gøre Indbegrebet af al
Realitet til Et med den højeste Fuldkommenhed. Dette fra den metafysiske
Side. Fra den psykologiske Side behandles Sagen i Afhandlingen om
Grundsætningernes Tydelighed, hvor det hedder (IV, 2): «Først i vore Dage
er man begyndt at indse, at Evnen til at indse det Sande er Erkendelse,
Evnen til at føle det Gode Følelse, og at de ikke maa forvexles med
hinanden. Ligesom der gives uopløselige Begreber om det Sande, d.\ v.\
s.\ om det, der træffes i Erkendelsens Genstande i og for sig betragtede,
saaledes gives der ogsaa en uopløselig Følelse af det Gode, hvilket aldrig
træffes i en Ting i og for sig, men stedse i Forhold til et følende
Væsen. Den praktiske Filosofi har sine Grundsætninger, lige som den
teoretiske har sine, og deres Udspring er endnu at undersøge, hvortil
Hutcheson og Andre have gjort en god Begyndelse.» --- I «Nachricht von
der Einrichtung seiner Vorlesungen» udtales det, at rigtige moralske
Domme kunne fældes af Menneskehjertet uden Omsvøb af Demonstrationer, ved
det, man kalder Sentiment. Der henvises til Shaftesbury, Hutcheson og
Hume som dem, der ere naaet videst i at opsøge Moralens første Grunde.
Det er Rousseau's Forgængere, Kant her henviser til. Om Rousseau selv
mindes vi, baade naar han som Grundlag for Etiken fordrer et Studium af
Menneskets Natur, der let miskendes ved de Forandringer, ydre Forhold
medføre, og naar den Opgave stilles: at afgøre, hvilken Fuldkommenhed der
tilkommer Mennesket i den naturlige Enfoldigheds Tilstand og i den vise
Enfoldigheds Tilstand.\footnote{Denne Modsætning har Kant langt senere,
med særlig Henvisning til Rousseau's forskellige Skrifter, behandlet i
sin Afhandling: \textit{Muthmasslicher Anfang des Menschengeschlechts}
(1786).} Men kun fordi vi ad anden Vej vide om Rousseau's Indflydelse paa
Kant, blive disse Spor tydelige.

Kant var, som det ses, selv paa en Vej, der maatte føre ham til at
sympatisere med Rousseau. Dennes Indflydelse har sandsynligvis lige som
Humes været vækkende, fremskyndende, har ført til raskere og bestemtere
Udfoldelse af Tanker, der allerede var i Skuddet, men har ikke bestaaet i
positiv Bibringelse af Noget, der ikke laa som Mulighed i den egne
forudgaaende Udvikling. Af ikke mindst Betydning har det været, at
Rousseau lærte, at Menneskene paa det intellektuelle saa vel som paa det
materielle Omraade have mange indbildte Fornødenheder. \textit{Nous
pouvons être hommes sans être savants!} --- \textit{Jeune homme, sachez
être ignorant!} --- siger den savoyiske Vikar til sin Discipel. Naar Kant
nu stod ved Bruddet med den dogmatiske Spekulation og havde omstyrtet
Grundlaget for den naturlige Teologi og den spiritualistiske Psykologi,
maatte han finde en Støtte i Rousseau's Tankegang. I Slutningsafsnittet
af «Träume eines Geistersehers» spores dette tydeligt. Efter sin Kritik
af de spekulative og de visionære Forsøg paa at faa Adgang til eller
Indblik i Aandeverdenen udbryder han: «Hvor Meget er der dog, som jeg
ikke forstaar!» --- men føjer strax til: «Og hvor Meget er der ikke, som
jeg ikke behøver!»

Etikens Emancipation fra Metafysik og Teologi var hermed given. Men Kant
anerkender ikke blot et af Teorien uafhængigt Grundlag for det Etiske.
Han gaar et Skridt videre, idet han udtaler, at det endog vilde være
meget betænkeligt, om den moralske Handlen blot motiveredes ved Troen paa
en anden Verden. I Stedet for at bygge Moralen paa Religionen bør man
omvendt bygge Religionen paa Moralen. De, som ville vide Besked med den
anden Verden, maa vente, til de komme derhen. Foreløbigt vide de, at de
skulle gøre deres Pligt, og mere behøve de ikke at vide. Kant ender sit
geniale Skrift som Voltaire sin Candide: Lad os gaa ud i Haven og
arbejde!

Ogsaa paa dette Punkt er der da Kontinuitet mellem de tidligere Arbejder
og de senere. Slutningsafsnittet af «Träume» udvikles videre tyve Aar
efter i «Kritik der praktischen Vernunft». Og at den vundne Indsigt i
Mellemtiden holdtes fast, ses af, at der i Dissertationen (§\ 9)
udtrykkeligt sondres mellem det højeste teoretiske og det højeste
praktiske Begreb.

19.\ Med Hensyn til det Grundlag, paa hvilket Kant byggede den af
Metafysik og Teologi uafhængige Etik, har der i Løbet af hans Udvikling
fundet betydelige Variationer Sted, dog saaledes, at der er et vist
Grundpræg, som er fælles for de forskellige Stadier, hans etiske
Opfattelse gennemløber.

I «Beobachtungen über das Schöne und Erhabene» (1764 [o:\ 1763]) gøres
allerede den sande Dyd afhængig af bestemte Grundsætninger, i Modsætning
til «den adopterede Dyd», der beror paa umiddelbare Naturtilbøjeligheder
som Medlidenhed og Velvillie. Det Karakteristiske for dette Kants første
etiske Standpunkt er, at de Grundsætninger, paa hvilke den sande Dyd
beror, ere «Bevidstheden om en Følelse, der lever i ethvert menneskeligt
Bryst, \ldots\ Følelsen af den menneskelige Naturs Skønhed og
Værdighed.»\footnote{\textit{Beobachtungen über das Gefühl des Schönen
und Erhabenen.} Königsberg 1764.\ p.\ 23.} Ved den principielle Fremhæven
af de almindelige Grundsætningers Betydning adskiller Kant sig fra de
engelske Etikere, til hvilke han ellers i disse Aar støtter sig. Men
Grundsætningerne selv antages at udvikle sig ved, at man bliver sig sin
Følelse af Menneskenaturens Skønhed og Værdighed bevidst. Denne Følelse
er altsaa den moralske Følelse. I den er, som det udvikles i «Träume
eines Geistersehers», den Enkeltes Villie bunden til den almindelige
Villie: «En hemmelig Magt nøder os til at rette vort Øjemed ogsaa mod
Andres Vel eller efter fremmed Villie, selv om det ofte sker ugerne og i
høj Grad strider mod den egennyttige Tilbøjelighed; --- og det Punkt,
hvor vore Drifters Retningslinier løbe sammen, ligger altsaa ikke blot i
os, men der er ogsaa i Andres Villie uden for os Kræfter, som bevæge os
\ldots\ Vi se os da i de hemmeligste Motiver afhængige af den almindelige
Villies Regel, og deraf udspringer i de tænkende Væseners Verden en
moralsk Enhed og en systematisk Forfatning efter rent aandelige
Love.»\footnote{\textit{Träume eines Geistersehers.} I, 2.\ (Kehrbachs
Udgave, p.\ 23.)}

Kants Etik staar her paa rent psykologisk Grund. Hvad der kundgør sig i
den enkelte Stemning og i det enkelte Tilfælde, underordnes det, som
fordres af Hensynet til den hele Verden, til hvilken Individet hører. Den
udvidede, ved den mest omskuende Erfaring betingede Følelse blive vi os
bevidst i almindelige Grundsætninger og kontrollere derefter Øjeblikkets
Rørelse og Trang. Ved denne Underordning af det Begrænsede under det
Universelle faar den etiske Følelse det Ophøjedes Karakter, og det er fra
den Side, Kant kommer ind paa at omtale den i «Beobachtungen».

20.\ Men Kant blev (uheldigvis) ikke staaende ved denne Forbindelse af
Psykologi og Etik. Hans Analyse førte ham paa det etiske saa vel som paa
det erkendelsesteoretiske Omraade til at sondre skarpt mellem Fornuft og
Sansning og mellem Form og Stof. Det synes endog, som denne Distinktion
paa det etiske Omraade er funden af ham før end paa det teoretiske
Omraade\footnote{Brev til M.\ Hertz af 21.\ Febr.\ 1772: «In der
Unterscheidung des Sinnlichen vom Intellektualen in der Moral und den
daraus entspringenden Grundsätzen hatte ich es schon vorher [o:\ før
Dissertationen af 1770] ziemlich weit gebracht.» Smlgn.\ Brev til Herder,
9.\ Maj 1767, til Lambert 2.\ Sept.\ 1770.}. Den psykologisk-genetiske
Opfattelse, som han i «Beobachtungen» og «Träume» endnu ansaa for
forenelig med Etiken, skød han til Side, for at det Etiske kunde faa sin
Kilde i den rene Fornuft. Indflydelsen af det Dilemma, han nu stiller op
mellem Fornuft og Sansning, bliver, at al Følelse regnes med til den
sanselige Del af vor Natur\footnote{«Alles Gefühl ist sinnlich.»
\textit{Kritik der praktischen Vernunft.} 1, 1, 3 (Kehrbachs Udgave, p.\
92).}, til Stoffet, det Empiriske, det, ved hvilket vi ere passive, ikke
selvvirksomme. I diametral Modsætning til Standpunktet i «Beobachtungen»
skulle nu Grundsætningerne bestemme Følelsen, ikke omvendt. Den moralske
Følelse bliver nu den ved Bevidstheden om den universelle Lov bestemte
Følelse af Agtelse. Agtelsen for andre Personligheder begrundes ved, at
de ere Organer for den moralske Lov, for den samme Lov, som vi selv mærke
i vort Indre. Og medens Kant i «Beobachtungen» nævnte Følelsen for
Menneskenaturens Skønhed ved Siden af Følelsen for dens Værdighed, saa er
denne Skønhedsfølelse senere helt falden bort i hans Etik. Denne Følelse
forudsatte en Harmoni mellem Elementerne i den menneskelige Natur, som
Kant formedelst sin skarpe Modsætning mellem Fornuft og Sansning ikke
længere kunde antage. Derfor var han saa betænkelig ved Schillers
Sammenstilling af «Anmuth und Würde». Skønt det var hans Mening, at
Pligten skulde øves med Glæde og Frejdighed, maatte der dog ingen
sanselige Elementer være medvirkende derved. Frejdigheden skal
udelukkende være betinget ved selve Udfoldelsen af Fornuftvilliens indre
Kraft, ikke medbetinget ved det harmoniske Forhold til den sanselige
Natur. I hvert Tilfælde maa Ynden ikke stilles foran Værdigheden; ti selv
om Pligtbevidstheden kan have Ynden til Ledsager, maa den dog slet ikke
tage Hensyn til denne. Den Stemning, Pligtbevidstheden vækker, har
Karakteren af det Ophøjede, ikke af det Skønne, og hvor det gælder etisk
Arbejde, maa Gratierne holde sig i ærbødig Afstand. Først naar Herkules
har betvunget Uhyrerne, kan han færdes blandt Muser og
Gratier.\footnote{Kants Svar paa Schillers «Anmuth u.\ Würde» findes i
hans «Religion innerhalb der Grenzen der blossen Vernunft». 2.\ Aufl.\
p.\ 10 f. --- Smlgn.\ ogsaa Optegnelser, der ere nedskrevne efter
Læsningen af Schillers Afhandling i «Thalia»: \textit{Lose Blätter} I.\
p.\ 120; 126 f.}

Saa skarp Modsætningen ogsaa er mellem Kants tidligere etiske Standpunkt
og det definitive, saa kunne vi dog ved Hjælp af de i de senere Aar
udgivne Kantiske Optegnelser efterspore Overgangen fra det ene Standpunkt
til det andet.

21.\ Naar Kant efter sit eget Udsagn lærte af Rousseau «at ære
Menneskene», og naar han (i «Beobachtungen») opfattede den moralske
Følelse som Følelsen af Menneskehedens Værdighed, hvad beroede saa i hans
Øjne denne Ære og Værdighed paa? --- Der fremtræder her i Kants etiske
Opfattelse tidligt et Hovedtræk, der aldrig forlod det, skønt han
filosofisk forklarede det meget forskelligt. Det var den Værdi, han
tillagde Evnen til at underordne det Enkelte og Skiftende under det
Almene og Uforanderlige, Indtryk og Stemninger under den trods alle
Tilskikkelser faste Grundsætning, det passivt Modtagne under den
aandelige Aktivitet. Dette Træk gjorde, at han allerede i «Beobachtungen»
henførte den moralske Følelse under Følelsen ved det Ophøjede, lige som
han ogsaa senere (se baade «Kritik der praktischen Vernunft» og «Kritik
der Urtheilskraft») fandt et nøje Slægtskab mellem disse to Følelser, dog
saaledes, at han ikke længere regnede den moralske Følelse ind under
Følelsen ved det Ophøjede, men snarere gik den omvendte Vej, idet han
tydede Følelsen ved det Ophøjede ved at antage en Medvirken af den
moralske Følelse.

Til Forstaaelse af Ændringen i Kants etiske Teori, der synes at være
foregaaet i Mellemrummet mellem 1766 (Träume) og 1770 (Dissertationen),
kan maaske følgende Betragtning tjene. --- Jo mere klart
Modsætningsforhold mellem det Enkelte, Skiftende, Passive paa den ene
Side, det Universelle, Faste og Aktive paa den anden Side betonedes, des
mere maatte der for Kant blive noget Utilfredsstillende ved at opfatte
Grundsætningerne som opstaaede ved, at vi blive os den moralske Følelse
bevidst. De synes derved at blive afhængige af Tilfældigheder. Et aldeles
sikkert Udgangspunkt for det Etiske vilde derimod faas, naar Fornuftens
Aktivitet sattes som det Første og Oprindelige, og Følelsen kun
betragtedes som dens Produkt. Altsaa Forholdet mellem Grundsætning og
Følelse maatte vendes om, naar den Retning, i hvilken Kants Etik allerede
gik i «Beobachtungen» og «Träume», skulde gennemføres absolut. Og at
denne Tendents maatte bestyrkes ved Kants ufortrødne Arbejde paa at
sondre mellem «det Rene» og «det Empiriske», følger af sig selv. Det var
maaske ikke en Gang frit for, at Begrebet om det Rene fik et moralsk Skær
i Kants Øjne. I det Mindste er der en Tvetydighed i den Maade, paa hvilken
han fra denne Tid af bruger Udtrykket «Sinnlichkeit» om alle Elementer i
Menneskenaturen med Undtagelse af Fornuften.

Han fjernede sig nu fra de engelske Moralfilosofer, til hvilke han
tidligere havde henvist, og betragtede det som en stor Fejl, at de førte
de moralske Kriterier tilbage til Følelsen af Lyst og Ulyst. Han saa ikke
nogen principiel Forskel mellem Shaftesbury's og Epikurs Standpunkt
(Dissertatio §\ 9), Noget, som allerede Mendelssohn (Brev til Kant 23.\
Decbr.\ 1770) med Rette ankede over. De moralske Begreber skulle nu ikke
erkendes ved Erfaring, men ved den rene Fornuft, og Moralfilosofien
hører, hvad de første Vurderingsprinciper (\textit{principia dijudicandi
prima}) angaar, til den rene Filosofi (Diss.\ §\ 7; 9). --- Hvorledes
Kant paa dette Stadium nærmere har tænkt sig Begrundelsen af det Etiske,
vides ikke. Dog foreligger der Udtalelser og Udkast, som vise, at der
forud for den etiske Teori, han udviklede i «Grundlegung zur Metaphysik
der Sitten» og i «Kritik der praktischen Vernunft», gaar en Teori, der
til en vis Grad danner et Overgangsled mellem den ældre psykologiserende
og den senere rationalistiske Etik.

22.\ I «Reflexionerne» findes fra et Tidspunkt, der ikke ligger langt fra
«Dissertationens» Tilblivelse, selv om det, efter Udgiverens Opfattelse,
ligger forud for denne, en Bestemmelse af Begrebet praktisk Idealisme, i
Følge hvilken denne bestaar i, at Lyksaligheden ikke afhænger af den ydre
Verden, eller udførligere: «at Verden kun er det, hvortil vi gøre den, at
den i et glad Sind har et lyst, i et sørgmodigt Sind et mørkt Udseende,
at Grundene til en lykkelig Tilstand maa søges i os selv.» Og denne
praktiske Idealisme kaldes udtrykkeligt «ikke Hjernespindets, men
Fornuftens, Visdommens Idealisme».\footnote{\textit{Reflexionen Kants.}
II.\ p.\ 319.\ (Nr.\ 1116; 1119.)}

Ikke lidet beslægtet med den saaledes bestemte praktiske Idealisme er den
Opfattelse, der kommer frem i et moralfilosofisk Udkast i de af R.\
Reicke udgivne «løse Blade af Kants Papirer». Her defineres Moraliteten
som «Ideen om Friheden som Principet for Lyksaligheden». «Det er sandt,
at Dyden har det Fortrin, at den vilde bringe den største Velfærd til
Veje ved Hjælp af det, Naturen frembyder. Men dens høje Værdi bestaar
ikke i, at den lige som tjener til Middel. At det er os selv, der
frembringe den som Ophavsmænd uden Hensyn til de empiriske Betingelser
(der kun kunne give partikulære Leveregler), --- at den fører
Selvtilfredshed med sig, --- det er dens indre Værdi». «Kun den er i
Stand til at være lykkelig, hvis Brug af hans Villie ikke strider mod de
Betingelser for Lyksalighed, som Naturen giver ham \ldots\ Lyksalighed er
egentligt ikke den største Sum af Lyst, men Lyst ved Bevidstheden om min
egen Magt til at være tilfreds (\textit{Selbstmacht zufrieden zu sein});
det er i det Mindste den vigtigste formale Betingelse for Lyksaligheden,
skønt der endnu \textit{materialiter} udfordres andre (som ved
Erfaringen)». «Lyksaligheden maa stamme fra en apriorisk Grund, som
Fornuften billiger». «Frihed er en Evne til at gøre og undlade
uafhængigt af empiriske Grunde \ldots\ Men jeg er kun fri for
Sanselighedens Tvang, ikke fra Fornuftens begrænsende Love \ldots\ Den
Ubundethed, med hvilken jeg kan ville, hvad der er selve min Villie imod,
saa at jeg ikke kan regne paa mig selv med Sikkerhed, maa i højeste Grad
vække Mishag hos mig. Det erkendes a priori, at der behøves en Lov, ved
hvilken Friheden begrænses saaledes, at Villien kan stemme overens med
sig selv. Denne Lov kan jeg ikke svigte uden at stride mod min Fornuft,
der ene kan fastsætte Villiens praktiske Enhed efter Principer». «Det er
Frihed efter Lovene for gennemgaaende Samstemmen med sig selv, der udgør
Personlighedens Værdi og Værdighed».\footnote{\textit{Lose Blätter} I.\
p.\ 10--15. Dette Fragment var af Udgiveren, Hr.\ Bibliothekar Rudolph
Reicke, henført til Aarene 80 eller 90. Da det af Indholdet var mig
klart, at det maatte skrive sig fra en Tid, der laa forud for
«Grundlegung zur Metaphysik der Sitten», som udkom 1785, henvendte jeg
mig til Udgiveren om Oplysning angaaende de Grunde, der havde bestemt
hans Angivelse af Manuskriptets Affattelsestid. Hr.\ Reicke har været saa
god at undersøge Sagen paa ny og har i et Privatbrev af 10.\ Novbr.\ 1892
underrettet mig om, at der Intet er i Vejen for, at Fragmentet kan skrive
sig fra 70'erne, men at han dog efter Skrifttrækkene snarere vilde
henføre det til 80'erne; til 90'erne anser han det nu for umuligt at
henføre det. Muligvis er det da blevet til i Slutningen af 70'erne eller i
Begyndelsen af 80'erne. Jeg vilde være tilbøjelig til at tro, at det var
forfattet nogen Tid før den sidste Redaktion af «Kritik der reinen
Vernunft», da det endnu ikke har den udprægede imperativiske Etik, som
allerede «Kritik d.\ r.\ V.» antyder, saa lidt som Læren om den
intelligible Karakter.}

Det etiske Standpunkt, der ligger til Grund for disse Udtalelser,
adskiller sig fra det, der fremtræder i «Beobachtungen» ved sin
rationalistiske og individualistiske Karakter. Det finder det Højeste i
den Aktivitet, ved hvilken den Enkelte kan være sin egen Lykkes Skaber,
og finder som den væsentlige Lov for denne Aktivitet, at Individet maa
handle med fuld Konsekvents, i Overensstemmelse med sig selv. Der er her
gjort et væsentligt Skridt i Retning af at lægge Hovedvægten paa det
Etiskes formale Side. Den etiske Lov skal se bort fra Indholdet, udspringe
af selve den frie Virksomhed, og kan da kun være dennes
Selvbegrænsning\footnote{Fragmentets Udtryk om Frihedens Begrænsning
kunde erindre om «Kritik der reinen Vernunft» 1.\ Udg.\ p.\ 569, hvor der
tales om «det Princip, ved hvilket Fornuften sætter den i sig selv
lovløse Frihed Skranker», og hvor den stoiske Vise omtales som et etisk
Ideal.} gennem Nødvendigheden af at bevare Overensstemmelsen med sig
selv. Kant havde her naat et Hovedpunkt i sin Etik: det indre Forhold
mellem Handling og Lov, et Forhold, som efter hans Opfattelse kun kunde
blive det rette, naar Loven var rent formal; kun da kunde den være
apriorisk, se bort fra alt bestemt Indhold. I genetisk Henseende er
derfor hint moralfilosofiske Fragment overordentligt interessant. Et
Problem danner det ved sin eudæmonistiske og individualistiske Karakter,
som stiller det i Modsætning baade til det tidligere og det senere
Standpunkt. Forholdet mellem Dyd og Lyksalighed, der for Kants senere
Etik blev saa vanskeligt et Problem, at kun religiøse Postulater kunde
klare det, indeholdt i Følge denne «praktiske Idealisme» ikke nogen
principiel Vanskelighed. De ydre Goder maa opgives, eller man maa i hvert
Tilfælde berede sig paa Resignation over for dem; men dette opvejes ved
den Friheds- og Magtfølelse, der her staar som det højeste Gode, lige som
den er bestemt ved hvad der her er den højeste Dyd. Standpunktet i
Fragmentet danner en Svingningsbue i Kants moralfilosofiske Udvikling af
analog Art som «Dissertationen» i hans erkendelsesteoretiske Udvikling.
Af størst Vigtighed var begge Steder den skarpe Distinktion mellem Form
og Indhold, a priori og a posteriori. Den blev bestemmende for Fremtiden.
Den skarpere Form, i hvilken denne Modsætning fremtræder i Kants
definitive Etik, medførte blandt Andet ogsaa den Modsætning mellem Dyd og
Lyksalighed, der kun kunde løses ad det religiøse Postulats Vej.

Hvis man vilde tænke sig, hvad der fra Fragmentets Standpunkt kan have
ført Kant over til hans definitive Etik, da ligger det nær at henvise
til, hvorledes han i sin teoretiske Filosofi var naat til at bestemme det
Almengyldige som det, der stemmer med alle fornuftige Subjekters Natur. I
Analogi hermed blev den moralske Lov ikke blot Udtryk for det enkelte
Individs Overensstemmelse med sig selv, men for Bevarelsen af
Fornuftsammenhængen i det Hele, for Individets Sammenhæng med hele den
intelligible Verden. Det hedder saaledes i en Optegnelse fra 70'erne:
«Den intelligible Verden har Love, efter hvilke jeg passer for enhver
Verden, ikke blot for denne, for Sanseverdenen, men ligegyldigt hvorledes
min Natur og den ydre Natur, samt det Samfund, hvori jeg lever, indrettes
\ldots\ Den passer med Klogskabens Regler, naar denne [o:\ Klogskaben]
udvides».\footnote{\textit{Reflexionen Kants.} II.\ p.\ 336.\ (Nr.\
1158.)} Disse sidste Ord synes at antyde en Udvikling fra det tidligere
rationalistiske Standpunkt til det definitive. Ved Hjælp af Ideen om den
intelligible Verden har Kant, trods Formalismen, opnaat at give sin
rationalistiske Etik en social Karakter, som hans første,
psykologisk-motiverede Etik havde ved selve sin Genesis.

Ogsaa en helt anden Tankerække har dog sandsynligvis medvirket ved denne
Overgang, og det en Tankerække, som netop havde en psykologisk-historisk
Karakter og for saa vidt kunde synes at stemme lidet med den
rationalistiske Vej, Kant var kommen ind paa. I Begyndelsen af 80'erne
var Kant stærkt beskæftiget med den Tanke, at de menneskelige Naturanlæg,
for saa vidt de angaa Brugen af Fornuften, kun komme til fuldstændig
Udvikling i Slægten, ikke i det enkelte Individ, da Fornuftens Udvikling
kræver langt længere Tid end det enkelte Individs Livsperiode. Maalet for
den menneskelige Udvikling kan derfor ikke være et individuelt, men maa
være et universelt Maal, et Maal for Alle, ikke blot et Maal for denne
eller hin Enkelte. Denne Tankegang ligger til Grund i Kants mærkelige
Afhandling «Idee zu einer allgemeinen Geschichte in weltbürgerlicher
Absicht», som udkom 1784, Aaret før «Grundlegung zur Metaphysik der
Sitten». I den anførte Tankegang maa man for en Del søge Forklaringen af,
at Kant nu ikke opfatter Moralloven som den formale Regel inden for det
enkelte Individs isolerede Stræben, men som den Regel, der bringer den
individuelle Stræben i Overensstemmelse med alle Fornuftvæseners Stræben.
Det er, ubevidst for Kant selv, en af psykologisk-historisk Erfaring
uddragen Ide, som har givet hans definitive Etik dens universalistiske
Karakter ved Siden af den rationalistiske Karakter. Selve den universelle
Formel, Kant gør til Moralprincip, indeholder lige som en Erindring om et
Samfund af Individer, hvis Villier det gælder om at bringe i Harmoni med
hverandre indbyrdes; den er egentlig talt en Idealisation af de Love for
Samkvem mellem personlige Væsener, som Erfaringen viser
os.\footnote{Smlgn.\ hvad jeg allerede har bemærket herom i min Afhandling
«Om Grundlaget for den humane Etik» (Kbhvn.\ 1876).\ p.\ 52.} Den store
Værdi af Kants Formel beror netop --- saa meget Kant selv ogsaa vilde
benegte det --- paa denne dens krypto-empiriske Oprindelse.

23.\ Endelig har ogsaa fornyet Undersøgelse og Analyse af den umiddelbare
etiske Følelse øvet sin store Indflydelse ved Tilblivelsen af Kants
definitive Etik. Lige som Analyse og Konstruktion kunde paavises
samvirkende ved Grundlæggelsen af Kants Erkendelsesteori, saaledes finde
vi ogsaa begge Metoder ved Grundlæggelsen af hans Etik. Det er netop
Eftervirkningen af denne umiddelbare Berøring med det Etiske i Praxis,
der giver Kants etiske Skrifter (især «Grundlegung zur Metaphysik der
Sitten») en saa frisk og kraftig Farve. I «Grundlegung» baner Kant sig
netop Vejen til den egentlige Moralfilosofi gennem Analyse af «den
almindelige sædelige Fornufterkendelse», --- saa lidet han ellers vil
lade psykologisk Erfaring have nogen som helst Indflydelse paa de etiske
Ideer. Ved denne fornyede Analyse fandt han da som et Hovedtræk ved det
Etiske, at den Enkelte føler sig bunden i sit Indre til en Lov, der paa
en Gang er Udtryk for hans eget inderste Væsen og indordner ham som Led
af et stort Fornuftrige. Det moralske Kriterium bliver nu det, om den
Grundsætning, Handlingen udtrykker, kan gøres til almindelig Grundsætning.
Denne Lære danner en Analogi til Læren om Fænomenernes lovmæssige
Sammenhæng som Udtryk for virkelig Erfaring.

Til Loven svarer Kraften. Fra den moralske Lov som et sig i al
Pligtbevidsthed kundgørende Faktum slutter Kant til den Fornuftens rene
Selvvirksomhed, som han kalder Frihed (i en af de flere Betydninger, i
hvilke han desværre bruger dette Ord). Ti umiddelbart kan Frihed,
Selvvirksomhed ikke opfattes, og lige saa lidt kan den udledes af
Erfaring. Den moralske Lov blive vi os umiddelbart bevidst, og den fører
til Antagelsen af Frihed som Evne til at handle uafhængigt af alle
sanselige Betingelser.\footnote{\textit{Kritik der praktischen Vernunft,}
§\ 6 Anm.\ (Kehrbachs Udg.\ p.\ 53). Kfr.\ \textit{Lose Blätter} I, p.\
120. --- Naar det enkelte Steder, f.\ Ex.\ Grundlegung p.\ 108, hedder,
at «det i Mennesket, som er ren Virksomhed, umiddelbart kommer til
Bevidsthed», maa efter Sammenhængen Ordet «umiddelbart» forstaas som
Modsætning til «durch Afficirung der Sinne».} Den moralske Lov er Udtryk
for Fornuftens Autonomi.\footnote{«Das moralische Gesetz ist das Gesetz
der Autonomie der reinen practischen Vernunft». \textit{Kritik der
praktischen Vernunft.} Kehrbachs Udg.\ p.\ 53.} Ved at lyde sit eget
Væsens Lov lyder den Enkelte tillige den aandelige Verdens almindelige
Lov.

Herved adskiller Kants Etik sig fra den tidligere rationalistiske Etik,
saaledes som den udformedes af «Intellektualisterne» i det 17.\
Aarhundrede og af Price. Naar disse grundede det Etiske paa Fornuft,
mente de dermed Tingenes objektive, fornuftige Sammenhæng, som Mennesket
opfatter ved sin Fornuft og derefter bøjer sig for. Fornuften som
menneskelig Evne fik her, som Fr.\ Jodl træffende udtrykker
sig\footnote{\textit{Geschichte der Ethik in der neueren Philosophie.}
I.\ München 1882.\ p.\ 221.}, en blot receptiv Rolle at udføre, nemlig at
tilegne sig den Fornuft, der antoges given i Verdensforholdene, og
anvende den i sin Handlen. Kants praktiske Fornuft er derimod, ligesom
hans teoretiske Fornuft, Udtryk for Menneskets eget inderste Væsen, et
Indbegreb af Love for Menneskets Aandsvirksomhed. Den dybsindige Tanke i
Kants definitive Etik er, at Mennesket ved at blive sig den moralske Lov
bevidst, bliver sig bevidst som Borger i to forskellige Verdener, der
mødes i hans Person: «Den moralske Skullen er Subjektets egen Villen som
Led i en intelligibel Verden og tænkes kun som Skullen, for saa vidt
Subjektet tillige betragter sig som Led af
Sanseverdenen».\footnote{\textit{Grundlegung zur Metaphysik der Sitten.}
3.\ Aufl.\ Riga 1792.\ p.\ 113.} I streng rationalistisk Form har Kant
her paa ny udviklet en Tanke, der i mere empirisk Form kom frem i
«Beobachtungen», naar der appelleredes til den menneskelige Naturs
Værdighed. Det er denne Værdighed, han mente at maatte søge en dybere
Begrundelse for end Psykologien kunde give. Ideen om Menneskets Værdighed
er det kontinuerlige Element i Kants Etik; hvad der varierer, er
Begrundelsen.

Trods sin ophøjede Karakter hviler Kants rationalistiske Etik dog, naar
den betragtes som Led i hans kritiske Filosofi, paa en stor Inkonsekvents.
Parallelen mellem den teoretiske og den praktiske Fornuft overholdes
ikke. Den Fjende, Kant bekæmper med Kritikens Vaaben, ser han som højst
forskellig paa de to Omraader. Paa det teoretiske Omraade er det
Spekulationen, den transscendente Brug af Begreberne, som skal bekæmpes.
«Kritik der reinen Vernunft» svarer til sit Navn, ti den er virkeligt en
Kritik af Fornuften. Men «Kritik der praktischen Vernunft» er slet ikke
nogen Kritik af Fornuftens Anvendelse paa etisk-praktisk Omraade,
saaledes som man af Titelen og af Analogien med hint Hovedværk skulde
formode. Og Kant indrømmer det egentlig selv. Opgaven er her, siger han
(Einleitung von der Idee einer Critik der praktischen Vernunft), «at
holde den empirisk betingede Fornuft borte fra den Anmasselse udelukkende
at ville afgive Villiens Bestemmelsesgrund».\footnote{Smlgn.\ nærmere
herom: A.\ Fouillée: \textit{Critique des systèmes de Morale
contemporains.} Paris 1883.\ p.\ 129--142.} Derimod drøfter han ikke den
Anmasselse af den rene Fornuft at ville være den udelukkende
Bestemmelsesgrund, men antager netop, at det Etiske ene bestaar heri. Han
indrømmer ikke, at der kunde være en kritisk Undersøgelse at anstille om
Fornuftens Berettigelse paa det praktiske Omraade, og det uagtet han selv
stærkt fremhæver det Uforklarlige i, at den rene Fornuft bestemmer et
«sanseligt» betinget Væsens Villie. --- Det hævner sig her, at det
psykologiske Grundlag er skudt til Side.

\section*{IV.\\ Det kopernikanske Princip.}

\noindent 24.\ Jeg vender tilbage til Kants teoretiske Filosofi for at
undersøge, hvorledes Vendepunktet i 1769, der medførte den egentlige
Grundlæggelse af den kritiske Filosofi, kan indføjes som Led i Kants
kontinuerlige Udviklingsgang, og hvorledes Konsekventserne efter Haanden
droges af dette Vendepunkt.

Hvad der var naaet i 1762, var Erkendelsen af, at de Begreber, med hvilke
man hidtil havde opereret i Filosofien, vare ubrugbare, ufærdige;
Analysen var ikke gennemført; om filosofisk Konstruktion kunde der først
i en fjern Fremtid blive Tale. Kant havde ikke opgivet Haabet om en
afsluttende teoretisk Filosofi, der gav Oplysning om Spørgsmaal, som laa
ud over Erfaringens Omraade. Men han vendte sig med Spot og Ironi mod
Dogmatikerne, der uanfægtede af alle Vanskeligheder byggede deres
Systemer og afgjorde alle Problemer. Hans Opmærksomhed var rettet mod
Metoden og mod Bestemmelsen af Erkendelsens Grænser, og han var allerede
begyndt at tale om en Kritik af Fornuften. For at finde Fornufterkendelsens
Grænser opererede han med Antinomiernes Metode, idet han i Modsigelsen
mellem to hver for sig begrundede Tankerækker saa Bevis paa, at Fornuften
havde vovet sig for langt ud. «Jeg søgte», siger
han\footnote{\textit{Reflexionen Kants.} II.\ p.\ 4 (Nr.\ 3).}, «Noget,
som var sikkert, om ikke i Henseende til Genstandene, saa dog i Henseende
til Erkendelsens Natur og Grænser. Jeg fandt efter Haanden, at mange af
de Sætninger, vi betragte som objektive, i Virkeligheden ere subjektive,
d.\ v.\ s.\ indeholde de Betingelser, under hvilke vi ene indse eller
begribe Genstandene.» Fra Bevidsthedens Indhold kom Kant altsaa til at
vende sin Opmærksomhed mod selve Bevidsthedsvirksomheden og gøre den til
Genstand for sin Analyse. Om Fremgangsmaaden ved denne Analyse er der
talt ovenfor (§\ 6).

Det er gennem en saadan Analyse, at han ikke blot kom til at skelne
mellem Form og Stof, men ogsaa mellem Anskuen og Tænken. Det hedder
saaledes i «Prolegomena» (p.\ 119): «Ved en Undersøgelse af den
menneskelige Erkendelses rene (intet Empirisk indeholdende) Elementer
lykkedes det mig først efter lang Eftertænkning med Sikkerhed at skelne
og udsondre Sansningens rene Elementarbegreber (Rum og Tid) fra
Forstandens Elementarbegreber.» --- Samtidigt opgik den Overbevisning hos
ham, at vi i hine Sansningens Elementarbegreber kun have Udtryk for de
Former, i hvilke vi efter vor Sanseevnes Natur opfatte Tingene, ikke for
Tingenes eget Væsen, og at ene derved matematisk Naturerkendelse er
mulig. Det var de vigtige Skridt, som gjordes 1769 og udvikledes i
Dissertationen af 1770.

25.\ Fra Aaret forud haves en lille interessant Afhandling af Kant: «Vom
ersten Grunde des Unterschiedes der Gegenden im Raume». Han forsvarer her
Newtons Opfattelse, i Følge hvilken det absolute Rum ligger til Grund for
enhver enkelt Rumsbestemmelse, over for Leibniz' Opfattelse af Rummet som
en Egenskab eller et Forhold, der er bestemt ved Tingenes Natur og
Virken. Han støtter sig især til, at ved aldeles lige og dog inkongruente
Legemer (som f.\ Ex.\ højre og venstre Haand) kan man ikke bestemme
Delenes Sted uden at tage Hensyn til Rummet uden for Legemerne. Deraf
slutter han, at Rummet som Helhed ligger til Grund ved Bestemmelsen af
det enkelte Sted.

I Løbet af det næste Aar (1768--69) maa det da være gaaet op for Kant, at
«det absolute Rum» ikke var noget Objektivt og Realt, men en subjektiv
Form, et Skema for vor Opfattelse af Tingene. Her have vi altsaa et
Exempel paa den Omsætten fra objektiv Realitet til subjektiv Betingelse,
som han omtaler i den ovenfor anførte Ytring. Det er i denne Sammenhæng
af Interesse at bemærke, at det samme Argument, som Kant i den lille
Afhandling af 1768 brugte til at bevise Newtons Rumsopfattelse over for
Leibniz's, bruger han senere (Prolegomena §\ 13)\footnote{Se ogsaa
\textit{Metaphysische Anfangsgründe der Naturwissenschaft.} Riga 1786.\
p.\ 7 f.} til at bevise sin egen Lære om Rummets Subjektivitet over for
den sædvanlige Antagelse af dets Objektivitet. Det er da tydeligt, at
Newtons Opfattelse har været ham en Station paa hans Vej. Det laa for
den, der i saa høj Grad som Kant havde beskæftiget sig med Studiet af
Newton, heller ikke saa fjernt at gøre Omsætningen fra det Objektive til
det Subjektive. Newton opfattede Rummet som \textit{sensorium dei}, en
Tanke, der ses at have beskæftiget Kant meget\footnote{Se herom Benno
Erdmann i «Reflexionen Kants». II.\ p.\ 104 f.\ (Noten).}. Det gjaldt da
blot om at sætte \textit{sensorium hominis} i Stedet for \textit{sensorium
dei}, hvad der for den, der var beskæftiget med at undersøge den
menneskelige Erkendelses almindelige Natur og Virkemaade, heller ikke
kunde ligge saa fjernt. --- Eftervirkningerne af Newtons Opfattelse
finder man i de psykologiske Vanskeligheder, i hvilke Kants Rumsteori
indvikler sig ved at beholde Bestemmelsen af Rummet som uendeligt og dog
at opfatte det som Form for den i Erfaringen virkende Sanseanskuelse,
hvis Genstand altid maa være endelig, --- Vanskeligheder, som Kant selv
kommer til at understrege ved at tale om Rummet som given
Uendelighed.\footnote{Se herom B.\ Erdmann i «Reflexionen Kants». II.\
p.\ 110 f.\ (Noten). --- Kant selv har i en mod Kästner rettet Afhandling,
der først fremdroges af Dilthey i \textit{Archiv für Gesch.\ der Phil.}
III (p.\ 83 ff.), søgt at klare dette Punkt.}

Den Bevisførelse, ved hvilken Kant søger at godtgøre Rum og Tid som
subjektive Anskuelsesformer, bestaar dels i direkte Analyse, dels i et
apagogisk Bevis. For Kortheds Skyld taler jeg i det Følgende kun om
Rummet, da Kants Ræsonnement for Tidens Vedkommende gaar aldeles parallelt
med hans Ræsonnement om Rummet.

Rummet er en Anskuelsesform, fordi det fremstiller for os en Helhed, der
bliver os bevidst ved successiv Sammenfatten af Delene, ved Syntese. Et
Forstandsbegreb derimod viser os blot Helhedens Sammensætning som betinget
ved Delene, uden at vise os dens Tilbliven. For Anskuelsen er den enkelte
Del betinget ved Helheden, for Forstanden omvendt (Diss.\ §\ 1; 15
Coroll.). --- Dersom paa den anden Side Rummet ikke var en
Anskuelsesform, men et Forstandsbegreb, saa vilde Begreberne om
Kontinuitet og Uendelighed indeslutte en Modsigelse, der falder bort,
naar man støtter dem paa Anskuelsens stadige og i Principet ustanselige
Fremskriden fra Del til Del (Diss.\ §\ 1; 2, III).

En dobbelt Bevisførelse --- en direkte og en apagogisk --- anvendes ogsaa
for at godtgøre, at Rummet kun er en subjektiv Opfattelsesmaade.

Dersom vi opfatte Rummet som et subjektivt Skema, i hvilket vi ordne de
Fornemmelser, der opstaa i os (\textit{schema coordinandi}), saa maa
Rummets Love ogsaa gælde for alle materielle Fænomener. Ti Alt hvad vi
skulle kunne opfatte, maa blive Genstand for den ordnende Aandsvirksomhed
(\textit{vis animi, omnes sensationes coordinans}). Sanseanskuelsens Love
maa derfor nødvendigvis være Naturens Love, og Geometriens Anvendelse i
Naturvidenskaben bliver derved forklaret: den samme Anskuelse, som
Geometrien udspringer af, bruge vi jo hver Gang vi opfatte Noget ved
Hjælp af den ydre Sans. Den anvendte Geometris Gyldighed er da en simpel
Slutning af Rumsformens Subjektivitet (Diss.\ §\ 15 C.\ og E.). --- Paa
den anden Side: «Hvis Rummets Begreb ikke var oprindeligt givet
formedelst Bevidsthedens Natur, saa vilde Geometriens Anvendelse i
Naturvidenskaben være usikker; ti der kunde da rejses Tvivl om, hvor vidt
selve dette fra Erfaringen hentede Begreb stemmede overens med Naturen»
(Diss.\ §\ 15 E.).

Kants Bevisførelse er af Interesse, da den viser os, hvorledes direkte
Analyse og antinomiske Slutningsrækker gaa Haand i Haand. Efter at han
ved Analyse har fundet, at Rummet er en syntetisk Anskuelsesform, viser
han, at en anden Antagelse vilde føre til Urimeligheder. Og idet han
derpaa gaar over til at drage de videre Slutninger af hvad han saaledes
har godtgjort, undersøger han de to Muligheder: at Rummet er subjektivt,
og at det ikke er subjektivt, og viser, at begges Konsekventser føre til
samme Resultat, nemlig at den anvendte Matematiks Mulighed er knyttet til
Anskuelsesformernes Subjektivitet. Vi finde her Afslutningen paa den
lange Række af Analyser og Antinomier, gennem hvilke Kant har arbejdet
sig frem til det afgørende Princip for hans Erkendelsesteori.

Dette Princip kunne vi, med Benyttelse af den bekendte Sammenligning i
Fortalen til 2.\ Udgave af «Kritik der reinen Vernunft», kalde det
kopernikanske Princip. Det udtaler, at vor Erkendelse af Tilværelsen er
bestemt ved en Erkendelsens Natur, og at vi, naar vi kende dennes Love,
kunne afgøre forud, hvilke almindelige Love Fænomenerne maa være
underlagte. --- Dette Princip gik --- hvad der ikke kan undre --- ikke
strax op for Kant i hele dets Rækkevidde. I Dissertationen anvendes det
kun paa Sanseanskuelsen, ikke paa Forstanden. Medens Sanseanskuelsen
altsaa kun viser os Tingene som Fænomener, d.\ v.\ s.\ saaledes som de
tage sig ud for vor Opfattelse, der iklæder dem vor sammenfattende
Anskuelses Former, saa skal Forstanden endnu i sine rene Begreber
(Substants, Aarsag, Mulighed, Virkelighed, Nødvendighed o.\ s.\ v.) kunne
føre til en Erkendelse af Tingene i sig selv, af Tilværelsens indre Væsen.
Det havde efter Kants eget Udsagn kostet ham lang Eftertanke at fastslaa
Forskellen mellem Anskuelse og Forstand, som den dogmatiske Filosofi
havde udvisket. Det var da intet Under, at han kunde mene, at disse to
forskellige Grene af vor Erkendelsesevne vare forskellige Vilkaar
underlagte, saa at Subjektiviteten og Begrænsningen til Fænomenerne kun
gjaldt for Anskuelsen, ikke for Forstanden. Den ærværdige, fra Platon
stammende Distinktion mellem \textit{Phænomena} og \textit{Noumena}, som
han nu ad sin egen Vej var ført til\footnote{Udtrykket «intelligibel
Verden» findes allerede i «Träume» (1.\ Th.\ 2.\ Hauptst.), hvor det
bruges om den «mystiske Filosofis» Verden af aandelige Substantser,
aabenbart med Henblik paa Leibniz's Monadeverden. Leibniz har selv brugt
Udtrykket \textit{mundus intelligibilis} om Monadeverdenen (Epistola ad
Hanschium. Op.\ phil.\ ed.\ Erdmann, p.\ 445 C.), eller, hvad der kommer
ud paa det Samme, om Formaalenes Verden (Animadv.\ in Principia Cartesiana.
Phil.\ Schr.\ ed.\ Gerhardt.\ IV, p.\ 391). --- B.\ Erdmann og R.\ Riedel
have vist, hvad der udførligt bekræftes ved «Reflexionerne», at Kants
intelligible Verden er Leibniz's Monadeverden med behørig Modifikation.},
maatte for ham ganske naturligt svare til Forskellen mellem Anskuelse og
Forstand. Der er dog intet Nyt i den Erkendelse af Noumena, som han ad
Tænkningens Vej endnu anser for mulig. Den bestaar væsentligt i den fra
«Beweisgrund» og tidligere Arbejder kendte Tanke om Tilværelsens
Enhedsgrund som nødvendig Betingelse for den lovmæssige Vexelvirkning i
Naturen. (Se ovenfor §\ 4.) Dissertationen hænger herved kontinuerligt
sammen med Kants tidligere Tankegang, samtidigt med, at den ved sit
kopernikanske Princip peger frem mod «Kritik der reinen Vernunft».

Friederich Paulsen, som henlægger Vækkelsen til 1769, opfatter
Dissertationen som en ved Humes Paavirkning hidført Reaktion mod de
foregaaende Aars skeptiske Tendentser. Jeg tror, at han stiller den i for
stor Modsætning til den nærmest foregaaende Periode. Kant havde aldrig
opgivet Haabet om et afsluttende Resultat ad Tænkningens Vej, skønt dets
Opfyldelse henlagdes til en ubestemt Fremtid. Han havde selv i sit mest
skeptiske Skrift erklæret sig forelsket i Metafysiken, og i et samtidigt
Brev udtalt, at det blot gjaldt om at afføre Metafysiken dens dogmatiske
Klædebon. Det er psykologisk forstaaeligt, at han efter Opdagelsen af
Forskellen mellem Anskuelse og Forstand og af det kopernikanske Princip i
dets Anvendelse paa Sanseerkendelsen kunde mene at være Maalet nær. Et
Øjeblik troede han at have naat den definitive Grænsebestemmelse, han i
flere Aar havde søgt efter. Dertil kommer, at den Erkendelse af Noumena,
som Kant paa dette Stadium anser for mulig, dog kun er symbolsk, idet
alle Anskuelsesdata mangle. Saa strengt Kant end skelner mellem Anskuelse
og Forstand, saa indskærper han dog allerede nu (Diss.\ §\ 10), at alt
Stof for vor Erkendelse faas ved Sansning. Og Forstandsbegreberne opfattes
ikke mindre end Anskuelsesformerne som Udtryk for den erkendende
Aandsvirksomheds Natur (Ibid.\ §\ 8). Dette er snarere en Begrænsning af
den noumenale Viden end en Udvidelse af den.

26.\ Hvad Kant manglede for at naa sit definitive Standpunkt var
Gennemførelsen af det kopernikanske Princip paa alle de første
Forudsætninger for vor Erkendelse, og dermed Udstrækningen af
Fænomenalismen til at gælde for al videnskabelig Erkendelse af
Tilværelsen. Det kunde synes, som om der ikke var langt frem, da det
første Skridt, det, der plejer at koste mest, allerede var gjort ved
Grundlæggelsen af den kritiske Lære om Rum og Tid. Men det kostede dog
Kant elleve Aars Arbejde at gøre Skridtet helt ud. Et Tegn paa Mandens
store Samvittighedsfuldhed, og tillige et Vidnesbyrd om, hvor højst
forskellig Sagen stiller sig for den midt i Problemerne stikkende Tænker
og for den tilbageskuende Historiker. Det kan allerede af praktiske
Grunde forstaas, at Kant krympede sig ved definitivt at udelukke
Erkendelsen fra Tingen i sig. Gennemførelsen af Fænomenalismen blev for
ham først mulig, da hans etiske Teori havde udviklet sig saaledes, at han
i det etiske Grundfaktum kunde se en Mulighed for Adgang til den Verden,
der nu blev spærret for den teoretiske Erkendelse. Det «Ubetingede», han
ikke længere kunde rumme i Erkendelsen, maatte han finde Plads til i
Troen; først saa ansaa han sin Opgave for løst, idet han havde fundet
«praktiske Data» for et Begreb, han ikke kunde opgive, men hvis teoretiske
Data han havde opløst\footnote{Se herom \textit{Kritik der reinen
Vernunft.} 2.\ Aufl.\ p.\ XIX--XXII, hvor disse «praktiske Data»
betragtes som Verifikation af «det kopernikanske Princip», ligesom
Newtons Tyngdelov gav en Verifikation af Kopernikus' Astronomi.}.
Hovedvanskeligheden har dog, som det ses af «Reflexionerne», været
teoretisk. Det drejede sig jo om at finde et Synspunkt for
Forstandsvirksomheden, ved hvilket det, der var paavist for
Sanseanskuelsen, ogsaa kunde komme til at gælde for den, trods den i
Dissertationen saa stærkt fremhævede Forskel.

Dog er den egentlige Kantiske Filosofi bleven til med Dissertationen.
Saaledes opfattede Kant det selv, naar han (i Brevene til Garve og
Mendelssohn) erklærede «Kritik der reinen Vernunft» for Frugten af tolv
Aars Arbejde, og naar han satte Aaret 1770 som Grænse for de Skrifter,
han ønskede optrykte. (Se ovenfor §\ 1.) Jeg kan ikke se, at den af Benno
Erdmann hævdede Antagelse, at Vækkelsen først har fundet Sted i Midten af
70'erne, er forenelig hermed. Det gjaldt jo i disse Aar om Gennemførelse
af et allerede fundet Princip, et Princip, som i og for sig forudsætter
et Brud med Dogmatismen. Kant kan ikke, da han 1783 skrev de bekendte
Linier i Fortalen til Prolegomena, der skulde volde Kantforskerne saa
meget Hovedbrud, have villet dele det Tidsrum i to skarpt adskilte Dele,
som han ellers betragtede som en Helhed. Den definitive Lukning af
Adgangen til Tingene i sig er ganske vist en epokegørende Begivenhed. Men
Udtrykket Vækkelse passer ikke paa den, hvorimod det godt passer paa en
Ændring i Maaden at forske paa, der medførte hin Lukning som sit sidste
Resultat.

27.\ I Dissertationen skelnes mellem en blot logisk Brug af Forstanden,
der bestaar i Sammenligning, Paavisning af Identitet og Forskel, og en
real Brug, der ikke karakteriseres nærmere, men som skal have absolut
Gyldighed, ikke blot Gyldighed for Fænomenerne, medens den logiske Brug
blot skal tjene til Ordning af Fænomenerne. Det Tankearbejde, der ligger
imellem Dissertationen og «Kritik der reinen Vernunft», gaar ud paa en
nærmere Undersøgelse af, hvorledes det forholder sig med denne «reale
Brug».\footnote{Udtrykket «real Brug af Fornuften», i Modsætning til den
blot logiske Brug, forekommer ogsaa i «Kritik der reinen Vernunft», i det
Mindste i Dialektiken. (2.\ Aufl.\ p.\ 355.)} Problemet stillede sig (som
det ses af Kants Breve til Lambert og Hertz) lige strax efter
Dissertationens Udgivelse. Naar de Forstandsbegreber eller «intellektuelle
Forestillinger», vi bruge til ad Tænkningens Vej at trænge ind i Tingenes
Væsen, udtrykke vor Aands Virkemaade, hvorledes kunne de da godtgøres at
have Gyldighed for selve Tilværelsen? Eller som det hedder i Brevet til
Hertz af 21.\ Februar 1772: «Naar saadanne intellektuelle Forestillinger
[Diss.\ §\ 8 nævnes: Begreberne Mulighed, Virkelighed, Nødvendighed,
Substants, Aarsag] bero paa vor indre Virksomhed, hvorfra kommer saa den
Overensstemmelse, som de skulle have med Genstande, der jo dog ikke
frembringes af dem, og hvorledes stemme den rene Fornufts Grundsætninger
om Genstandene overens med disse, naar denne Overensstemmelse ikke maa
søge Støtte i Erfaringen?»

Som Vejledning ved denne Undersøgelse havde Kant Opdagelsen fra
Dissertationen. Der kan, hvad de matematiske Grundforhold angaar, gives
en apriorisk Viden om den sanselige Natur, fordi de matematiske
Grundbegreber udtrykke Lovene for den syntetiske Konstruktion, der lægger
sig for Dagen i enhver Sanseanskuelse. Nu er det klart, at saa længe Kant
fastholdt sin tidligere saa stærkt hævdede Opfattelse af Tænkningen som
Analyse, maatte Forstandens reale Brug staa som en Gaade. Det førte i sin
Tid Kant et stort Skridt fremad, at han saa bestemt indskærpede, at al
Dømmen var en Sammenligning. (Afhandlingen om den falske Spidsfindighed.)
Nu gjaldt det om at vinde en dybere gaaende Opfattelse af
Forstandsvirksomheden, saa at denne kunde frembyde en Overensstemmelse
med Sanseanskuelsen. Lige frem at betragte Forstanden som en Anskuelse
vilde føre til Mystik. Hvilken Udvej var da at finde?

I «Reflexionerne» og i de «løse Blade» kan man se, hvor ihærdigt Kant har
arbejdet med de forskellige Grundbegreber, som udgjorde Indholdet af den
tidligere Ontologi, inden han fik denne saakaldte Videnskab forandret til
at være «Forstandens Analytik». Det gjaldt først at føre dem tilbage til
visse Hovedbegreber, for saa i disse at finde Udtryk for bestemte Arter
af Forstandsvirksomhed. Allerede i det anførte Brev til Hertz omtaler han
sin Bestræbelse for «at føre alle den absolut rene Fornufts Begreber
tilbage til et vist Antal af Kategorier, saaledes at deres Inddeling i
Klasser bestemmes ved nogle faa for Forstanden gældende Grundlove». Her
var to Veje at gaa. Man kunde gaa ud fra en Analyse af de Grundbegreber,
Forstanden opererer med, for at finde Lovene for Forstandens Virksomhed,
eller man kunde søge direkte at bestemme Loven for Forstandens Virksomhed
for saa derved igen at finde, hvilke Hovedarter af Grundbegreber den
efter sin Natur anvender paa det givne Stof. Kant er, som det synes,
gaaet begge Veje.

Først har han prøvet sig frem med forskellige Sammenstillinger og
Inddelinger af de «ontologiske» Begreber. Han har f.\ Ex.\ dvælet ved
Begreberne Sammenligning, Forbindelse og Forhold (eller Sammenhæng), og
han har særligt undersøgt Begreberne Forbindelse og Sammenligning og
fundet, at der ved begge de Operationer, de udtrykke, til Veje bringes
eller konstateres en Enhed: «Al Enhed er enten Sammenligningens eller
Forbindelsens Enhed. Den første have vi, for saa vidt Noget er ensartet
med noget Andet; den anden, for saa vidt en Mangfoldighed forbindes i en
Grund».\footnote{\textit{Reflexionen Kants.} II.\ p.\ 143 (Nr.\ 461).
Kfr.\ p.\ 164 (Nr.\ 525 ff.).} Han har undersøgt Begreberne Substants,
Aarsag og Vexelvirkning, som aabenbart særligt have interesseret ham. Og
ogsaa her var det Begrebet om en ved Tankeformen til Veje bragt Enhed i
det Mangfoldige, han særligt lagde Mærke til: «Begrebet Substants og
Accidents indeholder selv en Synthesis, ligeledes Aarsag og Virkning, og
Mangfoldighed i real Enhed [o:\ Vexelvirkning]. Efter sine forskellige
Forhold til den indre Sans maa nu Naturen nødvendigvis staa under en af
disse Synteser».\footnote{\textit{Lose Blätter.} I.\ p.\ 18 f.\ Smlgn.\
\textit{Reflexionen Kants.} II.\ p.\ 174--181 (Nr.\ 562--587).} Begrebet
Syntese synes han først blot at have anvendt paa disse tre Forhold, som
han tidligt har samlet under Udtrykket Relation. Det hedder saaledes i en
Optegnelse, der gaar forud for den egentlige Kategorilæres Tilblivelse:
«Kun for Relationen gælde objektivt syntetiske Sætninger om
Fænomenerne».\footnote{\textit{Lose Blätter.} I.\ p.\ 17. --- Udtrykket
Relation har Kant vistnok hentet fra den aristoteliske Kategoritavle. Han
har anvendt det paa den ovenanførte Kategorigruppe, før han brugte det
til Betegnelse for en Klasse af Domme. Se herom Adickes: \textit{Kants
Systematik als systembildender Faktor.} Berlin 1887.\ p.\ 26; 36. ---
Maaske er han ved Relationsbegrebet bleven ledet til at fæste sin
Opmærksomhed paa Dommene. Kfr.\ Refl.\ Nr.\ 532: «Vor Fornuft indeholder
ikke Andet end Relationer». Nr.\ 596: «Forholdets (Bevidsthedsenhedens)
Kategori er den vigtigste af alle. Ti Enhed angaar egentligt kun
Forholdet. Altsaa udgør dette Indholdet af Dommene over Hovedet og lader
sig ene tænke bestemt a priori».} Det kan (efter hvad der er udviklet i
første Kapitel af denne Afhandling) ikke undre os, at det særligt var
denne Klasse af Kategorier, der førte Kant til at udstrække Begrebet
Syntese til Forstandsvirksomheden og dermed tillige opdage, at det
kopernikanske Princip ogsaa maatte gælde for denne. --- Allerede i
Begyndelsen af Dissertationen havde han antydet som et fælles Synspunkt
for al vor Erkendelse, at den dels af givne Elementer danner et Hele,
dels opløser givne Helheder i deres Elementer. Derved opstaa to
Grænsebegreber: det absolut Enkelte (\textit{simplex} o:\ \textit{pars
quæ non est totum}) og Verdenstotaliteten (\textit{mundus} o:\
\textit{totum quod non est pars}). Her var altsaa henvist til de to
aandelige Grundoperationer, Syntese og Analyse, og det Spørgsmaal maatte
da frembyde sig, hvilken af dem der var den mest primitive. I
Dissertationen blev det nu slaaet fast for Anskuelsens Vedkommende, at
den beroede paa Syntese. Det store Fremskridt, som Kant gjorde i
Erkendelsens Psykologi i Løbet af 70'erne, var da Opdagelsen af, at ikke
blot al Anskuen, men ogsaa al Tænken bestaar i Syntese, naar man gaar
tilbage til Grundformen for Tankevirksomheden, og at Analysen altid er
sekundær i Forhold til den, da man kun kan opløse, hvad der allerede er
forbundet. Det gik op for Kant, at ikke blot Ordningen i Rum og Tid, men
ogsaa den indre Forbindelse mellem Fænomenerne, som vor Erkendelse til
Veje bringer, skyldes en «Sammenkædning» (\textit{Verkettung}) af
Forestillinger, en indre Bevidsthedshandling, som først bringer et
virkeligt Hele ud af de givne Elementer. Uden saadan Bevidsthedshandling
ingen Erfaring; altsaa maa den have
Almengyldighed.\footnote{\textit{Lose Blätter.} I.\ p.\ 16.}

Det kopernikanske Princip stod, saa vidt man kan se, fast for Kant i sin
fulde Udstrækning, saa snart han ved Undersøgelse af Kategorierne havde
fundet Syntesen som fælles Bestemmelse for al Erkendelsesvirksomhed. Vor
Aand, vort Jeg kom da til at staa som et Urbillede eller Forbillede for
alle Objekter, som vi skulle erkende. (Se ovenfor §\ 5 Slutn.; 12; 15.)

Men Undersøgelse af de gængse ontologiske Grundbegreber og Paavisning af
deres Udspring af en syntetisk Aandsvirksomhed var ikke den eneste Vej,
ad hvilken Kant kom til sin Kategorilære. Den var ham kun en foreløbig
Metode, der ikke tilfredsstillede ham, fordi den var for empirisk. Den
indeholdt ingen Garanti for, at alle Grundbegreberne vare fundne. Og for
Kant var det en Hovedsag, at Grænserne for Erkendelsen kunde afstikkes a
priori, efter fuldstændig Undersøgelse af alle Erkendelsens Grundformer.
Dette mente han at kunne opnaa ved at lægge Læren om Dommene til Grund,
saaledes at der til enhver særegen Art af Domme svarer en vis bestemt
Kategori. Men før han slog ind paa denne Vej, der som bekendt førte ham
til at belemre sin Fremstilling med et stort skolastisk Stillads, har
tydeligt nok den almindelige Ide om Forstandsvirksomheden som Syntese
staaet klart for ham. Dette ses af, hvad der ovenfor er anført af de
«løse Blade», og ogsaa «Reflexionerne» godtgøre det. Saaledes naar det
hedder (Refl.\ Nr.\ 600): «Dommen er Enheden af Bevidstheden om det
Mangfoldige i Forestillingen om et Objekt over Hovedet. Kategorien er
Forestillingen om et Objekt over Hovedet, for saa vidt det er bestemt i
Henseende til denne Bevidsthedens objektive Enhed». Her ligger Begrebet
Syntese til Grund for Definitionen af Dommen. Det maa være fundet som et
Grundbegreb, før det kunde ses, at al Dømmen falder ind under det.
Fremstillingen i «Kritik der reinen Vernunft» vidner om det Samme. Her
bestemmes Forstanden først som Evne til at erkende formedelst Begreber. En
Erkendelse ved Begreber er diskursiv, forudsætter altsaa en aandelig
Funktion, og en Funktion vil sige «Enheden af den Handling at ordne
forskellige Forestillinger under en fælles Forestilling» [smlgn.\
Dissertationens \textit{vis coordinandi}]. Af sine Begreber kan nu
Forstanden kun gøre Brug ved at dømme, d.\ v.\ s.\ bestemme Forestillingen
om en Genstand ved en anden Forestilling; det er kun Anskuelsen, der
umiddelbart har med Genstanden at gøre.\footnote{\textit{Kritik der reinen
Vernunft.} 1.\ Udg.\ p.\ 67--69 (Von dem logischen Verstandesgebrauche
überhaupt).} Forstandens Funktioner kunne fuldstændigt findes, naar man
fuldstændigt kan fremstille de forskellige Maader, paa hvilke der i
Dommene udtrykkes Enhed af forskellige Forestillinger. Og dette mener nu
Kant at kunne præstere ved en «forbedret» Udgave af den gængse logiske
Lære om Dommene.\footnote{Om den Maade, paa hvilken Kant til sin
«transscendentale» Brug har ændret Dommenes System se Adickes:
\textit{Kants Systematik,} p.\ 30--41.} Herpaa skal jeg ikke gaa nærmere
ind. Lige saa betydningsfuld som Kants almindelige Tanke om
Forstandsvirksomheden som Dømmen og om Dommen som Syntese er, lige saa
vilkaarlig er Anvendelsen i det Enkelte, skønt hans store Geni ogsaa her
viser sig i mange værdifulde Tanker, som han finder Plads til indenfor den
skolastiske Ramme. Hvis Kant i Stedet for at lægge saa megen Vægt paa at
deducere et fuldstændigt System af Kategorier og Grundsætninger havde
holdt sig til den almindelige Tanke, at de logiske Grundsætninger, der
betinge Dommenes formelle Gyldighed, ogsaa maa gælde for al Erfaring, for
al Sammenhæng mellem Fænomener, der skal blive til for os, saa havde han
undgaaet megen Uklarhed og Vilkaarlighed, og den transscendentale
Analytik som Læren om den anvendte Logiks Mulighed vilde da paa en
lærerig Maade være stillet ved Siden af den «transscendentale Æstetik»
som Læren om den anvendte Matematiks Mulighed.

28.\ Den «subjektive Deduktion» (d.\ v.\ s.\ den paa psykologisk Analyse
støttede Udledning af Erkendelsesevnens Natur), ved hvilken Kant naaede
til at bestemme Syntesen som Erkendelsens Grundform og til at give en
Oversigt over de specielle Former, i hvilke den lægger sig for Dagen, var
for ham dog ikke Hovedsagen. Hans Hovedopgave var derimod den objektive
eller transscendentale Deduktion, Paavisningen af Fornufterkendelsens
Gyldighed og Grænser. Hertil danner den analytiske Undersøgelse af
Formerne kun Indledning og Grundlag. Det var, som han skrev til Garve (d.\
7.\ August 1783), hans Opgave at grundlægge «en helt ny og hidtil ikke
forsøgt Videnskab, nemlig Kritiken af en a priori dømmende Fornuft», der
«af det blotte Begreb om en Erkendelsesevne (naar det er nøjagtigt
bestemt) a priori udleder alle Genstande, Alt hvad man kan vide om dem,
ja selv hvad man uvilkaarligt, skønt urigtigt nødes til at dømme om dem».
Med Rette klagede han over, at han paa dette Punkt blev saa grovt
misforstaaet i den første Bedømmelse af hans Hovedværk, og dette havde til
Følge, at den psykologiske Analyse traadte mere tilbage i hans senere
Fremstillinger. (Se ovenfor §\ 13.) Det Tiltrækkende ved den første
Udgave af «Kritik der reinen Vernunft» bestaar netop i den Friskhed og
Udførlighed, hvormed Syntesebegrebet fremstilles i dets Betydning for
Erkendelsesfunktionerne lige fra Sanseanskuelsen til den højeste
Tankeakt. Det er et Spørgsmaal, om denne Fremstilling, der i 2.\ Udgave
er forkortet betydeligt, ikke er det i Kants Lære, der er mest dybtgaaende
og har mest blivende Værdi. Som saa ofte er det ikke det bevidste
Hovedformaal, men de underordnede Formaal, som man maatte stille sig for
at naa hint, der har vist sig at være af størst Betydning.

Betydningen af Kants Syntesebegreb er for det Første den, at han derved
viser os, hvad det er at forstaa en Ting. Hans Erkendelseslære hviler paa
den simple Tanke, at det at forstaa Noget er at se det i saa nøje
Sammenhæng som muligt med alt Andet, vi kende. En forbindende
Aandsvirksomhed maa gøre sig gældende ved al Erkendelse; kun naar vi
kunne anvende den, have vi tilegnet os det givne Stof og føle den
ejendommelige Tilfredsstillelse, som følger med Afhjælpelsen af en Trang.
Hvad vi ikke forstaa, er det Isolerede og Sammenhængsløse. Naturligvis er
der mange Grader mellem de to Yderpunkter: forstaa og ikke forstaa,
Sammenhæng og Sammenhængsløshed, og, som vi have set (ovenfor §\ 14--15),
begik Kant den Fejl ikke at agte paa disse Gradsforskelle. Men den
Fortjeneste har han, først at have set Erkendelsesproblemet paa denne
simple Maade. Det er et Kolumbusæg, han har fundet og stillet paa
Spidsen. Mærkværdigt nok havde Hume fundet det samme Kolumbusæg, men lod
det trille fra sig. I sin «Treatise» (som Kant ikke kendte) kom Hume
egentligt til samme Grundbegreb, som Kant naaede til i «Kritik der reinen
Vernunft». Hume førte nemlig dér Erkendelsens Problem fra de specielle
Problemer (særligt Problemet om Geometriens reale Gyldighed og om
Aarsagsforholdet) tilbage til den Grundvanskelighed, som for ham ---
formedelst den atomistiske Psykologi, han gik ud fra --- laa i
Forbindelsen mellem Forestillinger i det Hele. Den store Gaade laa for
ham tilsidst i Enhedsprincipet (\textit{the uniting
principle}).\footnote{\textit{Treatise of human nature.} I, 3, 14 (ed.\
Selby-Bigge.\ Oxford 1888.\ p.\ 169). Kfr.\ Appendix (ibid.\ p.\ 636).}
Kant viser nu, at Forskellen mellem det Forstaaelige og det Gaadefulde
beror paa Forskellen mellem det, der kan sammenfattes i kontinuerlig
Sammenhæng, og det, hvor en saadan Funktion ikke er mulig. Selve den
forbindende Funktion kan man saa aabenbart kun finde gaadefuld, enten
naar man gaar ud fra, at det Sammenhængsløse skulde være det Forstaaelige,
eller naar man spørger om den større Sammenhæng, af hvilken vor
sammenfattende Funktion igen er betinget. Den første Vej fører ud i det
Meningsløse; den anden Vej fører til at stille os det sidste Problem,
Erkendelsesteorien og Psykologien i det Hele kunne staa over for, og som
betegner deres Grænse: nemlig om selve Erkendelsens, eller i det Hele
Bevidsthedens Opstaaen i Tilværelsen. Tanken er sig sit eget sidste
Problem.\footnote{Smlgn.\ min Psykologi.\ 3.\ Udg.\ p.\ 403.} --- Det
vilde sikkert have været af stor Betydning for Kants Udvikling, om han
havde kendt Humes Hovedværk paa et tidligt Stadium af sin Bane.
Sammenstødet mellem de to store Tænkere vilde da have angaaet et endnu
mere centralt Spørgsmaal, og Kant vilde maaske være bleven ledet til en
fuldkommen Gennemarbejdelse af det psykologiske Grundlag for sin
Erkendelsesteori, medens han paa den anden Side vilde have kunnet spare
sig et langt skolastisk Arbejde. Ægget vilde tidligere og lettere være
blevet sat paa Spidsen. Det var en for Filosofiens følgende Historie
skæbnesvanger Handling, Hume begik, da han, misfornøjet med den ringe
literære Lykke, hans Ungdomsværk gjorde og fortrædelig over de ubelejlige
Angreb, det var Genstand for fra teologisk Side, fornegtede dette sit
storartede Værk.\footnote{Eduard Grimm (\textit{Zur Geschichte des
Erkenntnissproblems.} Leipzig 1890.\ p.\ 584) ser i denne Fornegtelse en
Indrømmelse fra Humes Side af det Ensidige i «Treatise». Men at Motivet
afgjort var det ovenfor fremførte, ses tydeligt dels af Humes
Selvbiografi, dels af det Brev, i hvilket Hume opfordrede sin Forlægger
til at lade den Erklæring, i hvilken han fornegtede «Treatise», følge med
«Inquiry». \textit{Letters of David Hume to William Strahan.} Edited by
C.\ Birkbeck Hill.\ Oxford 1888.\ p.\ 288 f. --- Det var ogsaa af den
sidst omtalte Grund, at Hume ikke lod «Dialogerne» udkomme i sin Levetid.
Ibid.\ p.\ 330.}

Medens Kant ved sit ihærdige Tankearbejde fandt det Grundbegreb, der
kunde føre ud over den atomistiske Psykologi, som laa til Grund for
Empirismen, havde han ved det samme Grundbegreb opnaaet en Opfattelse af
Aandens Væsen, der førte ud over den spiritualistiske Dogmatismes
tilsnegne Begreber (smlgn.\ ovenfor §\ 8). Overfor Empirismen, der vilde
betragte Aandens Enhed som et blot Resultat af de mangfoldige Indtryk,
hævder han Enhedsvirksomheden som det aandelige Livs Grundpræg, der ikke
kan forklares ved ydre Paavirkning alene; over for Spiritualismen, der
vel har Øje for dette Grundpræg, men dogmatisk vil føre det tilbage til
en mystisk Substants bag Bevidstheden, hævder han (i sin Kritik af den
rationale Psykologi), at man fra Funktionen ikke har Ret til at slutte til
Substantsen, og at denne Slutning i hvert Tilfælde ikke kan føre til
virkelig Erkendelse.\footnote{Leibniz er den af de dogmatiske Filosofer,
der kommer Syntesebegrebet nærmest, som naar han (Monadologie §\ 14)
definerer «perception» saaledes: \textit{l'état passager qui enveloppe et
représente une multitude dans l'unité ou dans la substance simple.} ---
Leibniz bebrejdede Kartesianerne, at de ikke definerede hvad de forstod
ved Tænken (o:\ Bevidsthed). Hans egen Definition var, at Bevidsthed er
Mangfoldighed udtrykt i en Enhed. (\textit{Philos.\ Schr.} herausg.\ von
Gerhardt.\ VII.\ p.\ 551). «\textit{Les âmes sont des unités et les corps
sont des multitudes, mais les unités ne laissent de representer les
multitudes \ldots\ C'est dans cette réunion que consiste la nature
admirable du sentiment.}» (ibid.\ p.\ 542.)} Kants Syntesebegreb
udtrykker Grundform og Grundlov for det aandelige Liv, saa vidt vi kunne
kende det ad Erfaringens Vej. Han har derved givet Psykologien et vigtigt
Synspunkt, en Maalestok ved Vurderingen af det aandelige Livs Udvikling,
--- og tillige nogle af dens vigtigste Problemer, ti den aandelige Enhed
er ofte truet ved Elementernes Mangfoldighed og Modstrid. Særligt
interessant er det at se, hvor højt Kant ved dette Begreb hæver sig over
Oplysningstidens Fremhæven af det klart Bevidste og forstandsmæssigt
Gennemskuede. Ti Syntesen ligger vel til Grund for Bevidstheden, men
behøver ikke at mærkes af Bevidstheden selv. Den er Bevidsthedens stadige
Betingelse, men ikke nødvendigvis dens Genstand. Den virker blindt,
instinktmæssigt i vort inderste Væsen. Naar Kant i «Kritik der
Urtheilskraft» (§\ 46) definerer Geni som «den medfødte aandelige Evne,
ved hvilken Naturen giver Regler for Kunsten», saa er det kun en speciel
Anvendelse af Noget, som gælder for alt Bevidsthedsliv. Der er, om man
vil, noget Genialt i ethvert bevidst Væsen, som hos dem, der særligt ere
at kalde Genier, kun træder frem i højeste Potents. --- Det Umiddelbare
og Uvilkaarlige i det aandelige Liv bliver saaledes ved den Kantiske
Opfattelse hævdet i sin Ret og Betydning som det stadige Grundlag. Kants
Syntesebegreb har en langt mere omfattende Betydning end den han tillægger
det i sin Erkendelsesteori, og det har i hvert Tilfælde mere blivende
Værd end andre Bestanddele af hans Filosofi, som han selv satte nok saa
højt.

\end{document}
